% !TEX program = xelatex
% C/C++ 历史与 ISO/IEC 治理模式
\documentclass[12pt,a4paper]{ctexbook}

% ===== 页面 =====
\usepackage[margin=2.5cm]{geometry}

% ===== 字体 =====
\usepackage{fontspec}
\usepackage{iffont}
\IfFontExistsTF{Consolas}
  {\setmonofont{Consolas}}
  {\setmonofont{DejaVu Sans Mono}[Scale=MatchLowercase]}

% ===== 颜色 =====
\usepackage[dvipsnames,svgnames,x11names]{xcolor}
\definecolor{codebg}{HTML}{F5F5F5}
\definecolor{codeframe}{HTML}{E0E0E0}
\definecolor{linenumgray}{HTML}{999999}
\definecolor{cppkeyword}{HTML}{1A1AFF}
\definecolor{cppcomment}{HTML}{2E8B2E}
\definecolor{cppstring}{HTML}{B22222}
\definecolor{cppheadbg}{HTML}{2E4057}
\definecolor{cheadbg}{HTML}{B36B00}
\definecolor{govheadbg}{HTML}{1F4E79}
\definecolor{yearheadbg}{HTML}{8B4513}
\definecolor{csharptype}{HTML}{006400}
\definecolor{algheadbg}{HTML}{1A3C5E}

% ===== listings =====
\usepackage{listings}

% ---- C++ ----
\lstdefinestyle{cppstyle}{%
  language=C++,
  basicstyle=\small\ttfamily,numbers=left,numberstyle=\tiny\color{linenumgray},
  frame=single,rulecolor=\color{codeframe},framesep=6pt,
  breaklines=true,tabsize=4,columns=flexible,keepspaces=true,
  backgroundcolor=\color{codebg},showstringspaces=false,
  keywordstyle=\color{cppkeyword}\bfseries,
  commentstyle=\color{cppcomment}\itshape,
  stringstyle=\color{cppstring},
  morekeywords={%
    override,final,noexcept,constexpr,consteval,constinit,
    decltype,auto,concept,requires,co_await,co_yield,co_return,
    nullptr,sizeof...,static_assert,thread_local,alignas,alignof,
    namespace,using,template,typename,virtual,explicit,friend,
    mutable,volatile,register,extern,inline,
    module,import,export,
    std,string,vector,array,map,unordered_map,set,unordered_set,
    unique_ptr,shared_ptr,weak_ptr,optional,variant,any,
    tuple,pair,span,string_view,
    cin,cout,cerr,endl,print,println,
    move,forward,swap,make_unique,make_shared,
    begin,end,size,empty,push_back,emplace_back,
    is_same_v,integral_constant,
    chrono,literal,operator
  },
  deletekeywords={int,char,float,double,bool,void,long,short,unsigned,signed},
  emph={%
    int8_t,int16_t,int32_t,int64_t,uint8_t,uint16_t,uint32_t,uint64_t,
    size_t,ptrdiff_t,intptr_t,uintptr_t,wchar_t,char8_t
  },
  emphstyle=\color{csharptype},
  escapeinside={(*@}{@*)}
}

% ---- C ----
\lstdefinestyle{cstyle}{%
  language=C,
  basicstyle=\small\ttfamily,numbers=left,numberstyle=\tiny\color{linenumgray},
  frame=single,rulecolor=\color{codeframe},framesep=6pt,
  breaklines=true,tabsize=4,columns=flexible,keepspaces=true,
  backgroundcolor=\color{codebg},showstringspaces=false,
  keywordstyle=\color{cheadbg}\bfseries,
  commentstyle=\color{cppcomment}\itshape,
  stringstyle=\color{cppstring},
  morekeywords={%
    _Bool,_Complex,_Imaginary,_Atomic,_Alignas,_Alignof,_Generic,_Noreturn,_Static_assert,_Thread_local,
    bool,true,false,nullptr,typeof,typeof_unqual,
    restrict,inline,register,auto,extern,static,const,volatile,signed,unsigned,
    struct,union,enum,typedef,sizeof,void,char,short,int,long,float,double,
    int8_t,int16_t,int32_t,int64_t,uint8_t,uint16_t,uint32_t,uint64_t,size_t,ptrdiff_t,
    alignas,alignof,static_assert,thread_local,constexpr,
    _Pragma
  },
  emph={int8_t,int16_t,int32_t,int64_t,uint8_t,uint16_t,uint32_t,uint64_t,size_t},
  emphstyle=\color{csharptype},
  escapeinside={(*@}{@*)}
}

\lstset{style=cppstyle}

% ===== tcolorbox =====
\usepackage[most]{tcolorbox}

% ===== 代码标签 =====
\newcommand{\cpplabel}{%
  \tcbox[on line,colback=cppheadbg,colframe=cppheadbg,
    boxrule=0pt,arc=1pt,left=2pt,right=2pt,top=1pt,bottom=1pt,
    colupper=white,fontupper=\scriptsize\bfseries]{C++}}

\newcommand{\clabel}{%
  \tcbox[on line,colback=cheadbg,colframe=cheadbg,
    boxrule=0pt,arc=1pt,left=2pt,right=2pt,top=1pt,bottom=1pt,
    colupper=white,fontupper=\scriptsize\bfseries]{C}}

\newcommand{\cppxxlabel}[1]{%
  \tcbox[on line,colback=cppheadbg,colframe=cppheadbg,
    boxrule=0pt,arc=1pt,left=2pt,right=2pt,top=1pt,bottom=1pt,
    colupper=white,fontupper=\scriptsize\bfseries]{#1}}

\newcommand{\govlabel}{%
  \tcbox[on line,colback=govheadbg,colframe=govheadbg,
    boxrule=0pt,arc=1pt,left=2pt,right=2pt,top=1pt,bottom=1pt,
    colupper=white,fontupper=\scriptsize\bfseries]{治理}}

\newcommand{\yearlabel}[1]{%
  \tcbox[on line,colback=yearheadbg,colframe=yearheadbg,
    boxrule=0pt,arc=1pt,left=3pt,right=3pt,top=1pt,bottom=1pt,
    colupper=white,fontupper=\scriptsize\bfseries]{#1}}

% ===== tcolorbox 提示框 =====
\newtcolorbox{guideline}[1][]{
  enhanced,breakable,
  colback=blue!5!white,colframe=blue!75!black,
  fonttitle=\bfseries,title=要点,#1,
  boxed title style={colback=blue!75!black},
  attach boxed title to top left={yshift=-2mm,xshift=2mm},
  arc=2mm,boxrule=0.8mm,
}
\newtcolorbox{compare}[1][]{
  enhanced,breakable,
  colback=green!3!white,colframe=green!50!black,
  fonttitle=\bfseries,title=对比,#1,
  boxed title style={colback=green!50!black},
  attach boxed title to top left={yshift=-2mm,xshift=2mm},
  arc=2mm,boxrule=0.8mm,
}
\newtcolorbox{warning}[1][]{
  enhanced,breakable,
  colback=red!5!white,colframe=red!75!black,
  fonttitle=\bfseries,title=常见陷阱,#1,
  boxed title style={colback=red!75!black},
  attach boxed title to top left={yshift=-2mm,xshift=2mm},
  arc=2mm,boxrule=0.8mm,
}
\newtcolorbox{note}[1][]{
  enhanced,breakable,
  colback=orange!5!white,colframe=orange!80!black,
  fonttitle=\bfseries,title=注意,#1,
  boxed title style={colback=orange!80!black},
  attach boxed title to top left={yshift=-2mm,xshift=2mm},
  arc=2mm,boxrule=0.8mm,
}
\newtcolorbox{bestpractice}[1][]{
  enhanced,breakable,
  colback=purple!5!white,colframe=purple!60!black,
  fonttitle=\bfseries,title=最佳实践,#1,
  boxed title style={colback=purple!60!black},
  attach boxed title to top left={yshift=-2mm,xshift=2mm},
  arc=2mm,boxrule=0.8mm,
}
\newtcolorbox{stdbox}[1][]{
  enhanced,breakable,
  colback=gray!5!white,colframe=gray!70!black,
  fonttitle=\bfseries,title=标准条款,#1,
  boxed title style={colback=gray!70!black},
  attach boxed title to top left={yshift=-2mm,xshift=2mm},
  arc=2mm,boxrule=0.8mm,
}
\newtcolorbox{govbox}[1][]{
  enhanced,breakable,
  colback=cyan!5!white,colframe=cyan!60!black,
  fonttitle=\bfseries,title=治理观察,#1,
  boxed title style={colback=cyan!60!black},
  attach boxed title to top left={yshift=-2mm,xshift=2mm},
  arc=2mm,boxrule=0.8mm,
}
\newtcolorbox{historybox}[1][]{
  enhanced,breakable,
  colback=brown!5!white,colframe=brown!60!black,
  fonttitle=\bfseries,title=历史快照,#1,
  boxed title style={colback=brown!60!black},
  attach boxed title to top left={yshift=-2mm,xshift=2mm},
  arc=2mm,boxrule=0.8mm,
}
\newtcolorbox{timelinebox}[1][]{
  enhanced,breakable,
  colback=gray!3!white,colframe=gray!60!black,
  fonttitle=\bfseries,title=时间线,#1,
  boxed title style={colback=gray!60!black},
  attach boxed title to top left={yshift=-2mm,xshift=2mm},
  arc=2mm,boxrule=0.6mm,
}

% ===== 章节标题 =====
\usepackage{titlesec}
\usepackage{fancyhdr}
\usepackage{graphicx}
\usepackage{amssymb}
\usepackage{amsmath}
\usepackage{tabularx}
\usepackage{booktabs}
\usepackage{longtable}
\usepackage{array}
\usepackage{multirow}
\usepackage{enumitem}
\usepackage{seqsplit}
\usepackage{float}

\titleformat{\chapter}[display]
  {\normalfont\huge\bfseries\color{algheadbg}}
  {\chaptertitlename\ \thechapter}{20pt}{\Huge}
\titleformat{\section}
  {\normalfont\Large\bfseries\color{blue!70!black}}
  {\thesection}{1em}{}
\titleformat{\subsection}
  {\normalfont\large\bfseries\color{blue!60!black}}
  {\thesubsection}{1em}{}

\pagestyle{fancy}
\fancyhf{}
\fancyhead[LE,RO]{\thepage}
\fancyhead[RE]{\leftmark}
\fancyhead[LO]{\rightmark}

\setlist[itemize]{leftmargin=2em,itemsep=2pt}
\setlist[enumerate]{leftmargin=2em,itemsep=2pt}

% ===== 超链接 =====
\usepackage{hyperref}
\hypersetup{
  colorlinks=true,
  linkcolor=blue!70!black,
  urlcolor=blue!60!black,
  bookmarks=true,
  pdfauthor={hsmbanlance},
  pdftitle={C/C++ 历史与 ISO/IEC 治理模式},
}

% ===== 自定义命令 =====
\newcommand{\cpp}[1]{C++#1}
\newcommand{\kw}[1]{\texttt{#1}}
\newcommand{\seqkw}[1]{\seqsplit{\texttt{#1}}}
\newcommand{\paper}[1]{\texttt{#1}}
\newcommand{\yrlbl}[1]{\yearlabel{#1}}

% ====================================================================
\begin{document}

\frontmatter

\begin{titlepage}
\centering
\vspace*{4cm}
{\Huge\bfseries C/C++ 历史与 ISO/IEC 治理模式\par}
\vspace{1cm}
{\LARGE 从 Bell Labs 到 ISO 委员会：去中心化标准化的优势与代价\par}
\vfill
{\large \today\par}
\end{titlepage}

\tableofcontents

\mainmatter

% ====================================================================
\part{C 语言的诞生与标准化}
% ====================================================================

\chapter{C 的前史：从 Multics 到 UNIX}

\section{时代背景：1960 年代的分时系统}

\yrlbl{1965} 1960 年代中期，麻省理工学院（MIT）、贝尔实验室（Bell Labs）与通用电气（GE）联合启动了 \textbf{Multics} 项目——一个雄心勃勃的分时操作系统，目标是让一台大型机同时为数百名用户提供交互式计算服务。Multics 的设计文档由 Fernando Corbató 等人起草，采用当时最先进的 PL/I 语言实现，运行在 GE-645 主机上。

\yrlbl{1969} 然而 Multics 的复杂度迅速失控。PL/I 编译器迟迟不能稳定，硬件资源紧张，项目进度严重落后。1969 年，Bell Labs 决定退出 Multics 项目。退出之后，Ken Thompson、Dennis Ritchie 与 Doug McIlroy 等几位原 Multics 团队成员，仍然怀念那种\textbf{多用户、多任务、文件系统统一}的工作模式，但他们手上只剩一台几乎被废弃的 \textbf{PDP-7} 小型机。

\begin{historybox}
\textbf{Multics 的遗产}

虽然 Multics 项目本身在商业上不算成功（直到 1975 年才有第一个商用部署，2000 年才被彻底退役），但它对操作系统设计的影响极其深远。Multics 引入的概念包括：分层文件系统、进程作为基本执行单元、动态链接、安全环（security rings）、shell 作为命令行接口。所有这些概念都通过 UNIX 流传至今。C 语言的诞生与 Multics 项目的失败直接相关——没有 Multics 的复杂度教训，就没有 UNIX 的极简哲学，也就没有为 UNIX 而生的 C。
\end{historybox}

\section{BCPL、B 语言与"New B"}

\yrlbl{1967} 在 Multics 之前，剑桥大学的 Martin Richards 于 1967 年设计了 \textbf{BCPL}（Basic Combined Programming Language），这是一种为系统编程设计的语言。BCPL 的核心特征是\textbf{无类型}——所有数据都被当作机器字（machine word）处理，由程序员自己负责解释。这种设计极大简化了编译器实现，让 BCPL 能够快速移植到不同机器。

\yrlbl{1969} Ken Thompson 在 PDP-7 上工作时，需要一个比汇编更易移植、但比 BCPL 更紧凑的语言。他在 BCPL 的基础上做了删减，去掉了不需要的部分（比如复杂的 I/O 库），发明了 \textbf{B 语言}。B 是一种解释执行的无类型语言，所有值都视为 16 位的字。Thompson 用 B 重写了 PDP-7 上的部分系统工具，包括早期版本的 UNIX 内核。

\yrlbl{1970} 但 B 的无类型设计在 PDP-11 上遇到了麻烦。PDP-11 引入了\textbf{字节寻址}（byte-addressable memory）和新的数据类型（如 \kw{char}），B 无法自然表达。Dennis Ritchie 接手扩展 B，加入类型概念和结构体（\kw{struct}），称为 \textbf{New B}（NB）。New B 是 C 的直接前身。

\clabel\ \begin{lstlisting}[style=cstyle]
/* B 语言示例：所有变量都是机器字，无类型 */
main() {
    auto a, b, c;
    a = 1; b = 2;
    c = a + b;     /* c 是 3，但解释为整数还是指针由上下文决定 */
    print(c);
}

/* New B / 早期 C 示例：引入了类型 */
int main() {
    int a, b, c;
    char ch;
    a = 1; b = 2;
    c = a + b;     /* 明确是整数加法 */
    ch = 'A';      /* 明确是字符 */
    return 0;
}
\end{lstlisting}

\section{C 语言的诞生：1972-1973}

\yrlbl{1972} Ritchie 在 New B 的基础上继续演进，加入了完整的类型系统、结构体（\kw{struct}）、联合（\kw{union}）、指针算术与自增/自减运算符。这门新语言最初被称为"New B"或"C"，1972 年前后正式定名为 \textbf{C}——只是 B 的下一个字母，没有任何特殊含义。

\yrlbl{1973} 关键的一年。Ritchie 用 C 重写了 B 编译器，Thompson 与 Ritchie 用 C 重写了 UNIX 内核的大部分代码。在此之前，操作系统几乎都是用汇编写的——每个硬件平台都需要重写一遍。用 C 写操作系统意味着\textbf{只要 C 编译器能移植到新机器，整个操作系统就能跟着移植}。这是软件史上第一次大规模验证"高级语言写系统软件"的可行性。

\begin{guideline}
\textbf{C 的核心设计哲学（1973 年确立）}

(1) \textbf{贴近硬件}：C 提供指针、位运算、直接的内存布局控制，但通过类型系统提供最低限度的安全检查。

(2) \textbf{小而精}：C 的关键字仅 32 个（C89），没有内置的字符串、I/O、容器、面向对象——这些都靠标准库或程序员自己实现。

(3) \textbf{可移植}：编译器只需为每个目标机器实现一次，所有 C 程序随之可移植。这与"操作系统用汇编写、每个机器重写一次"的传统模式形成鲜明对比。

(4) \textbf{信任程序员}：C 不做运行期边界检查、不做空指针检查、不做整数溢出检查。程序员被认为有能力管理自己的内存与类型。
\end{guideline}

\yrlbl{1973} 1973 年，Ritchie 与 Thompson 在 ACM 杂志上发表了《The UNIX Time-Sharing System》——UNIX 的第一篇公开论文。论文中提到了实现语言是"一种新的、尚未广泛使用的语言 C"，这是 C 第一次出现在公开文献中。

\section{早期移植与扩散（1974-1978）}

\yrlbl{1974} UNIX 第四版（V4）发布，是第一个用 C 写的 UNIX 版本。Bell Labs 内部的 PDP-11 团队开始使用。

\yrlbl{1975} UNIX V6 发布，附带完整源代码（以极低的许可费给大学使用）。加州大学伯克利分校（UC Berkeley）的学生开始在此基础上做扩展，逐渐演变为 \textbf{BSD}（Berkeley Software Distribution）。

\yrlbl{1977} Ritchie 完成了 C 编译器的第一个真正的\textbf{跨架构移植}：把 C 编译到 Interdata 8/32——一种与 PDP-11 完全不同的非 DEC 机器。这次移植暴露了 C 早期定义中关于整数大小、字节序、对齐的模糊地带，促使 Ritchie 开始考虑写一份正式的语言参考。

\yrlbl{1978} 同年，C 编译器已经移植到 IBM System/370、Honeywell 6000、Interdata 8/32、VAX-11 等多种机器。Bell Labs 内部已有数十个项目使用 C，外部用户也在增长。C 不再是"Bell Labs 的内部语言"，而是一个事实上的多平台系统编程语言。

\begin{timelinebox}
\textbf{C 语言诞生时间线（1967-1978）}

\yrlbl{1967} BCPL 由 Martin Richards 在剑桥设计。

\yrlbl{1969} Bell Labs 退出 Multics；Ken Thompson 在 PDP-7 上发明 B 语言。

\yrlbl{1970} Dennis Ritchie 扩展 B 为 New B（NB），加入类型概念。

\yrlbl{1971} UNIX 第一版手册（PDP-11）发布，使用汇编+B 实现。

\yrlbl{1972} Ritchie 在 NB 基础上完成 C 语言的雏形。

\yrlbl{1973} C 重写 UNIX 内核，证明"高级语言写 OS"可行。

\yrlbl{1974} UNIX V4 发布，C 实现版本首次内部使用。

\yrlbl{1975} UNIX V6 源代码对大学开放，BSD 项目萌芽。

\yrlbl{1977} Ritchie 完成 C 编译器到 Interdata 8/32 的跨架构移植。

\yrlbl{1978} 《The C Programming Language》（K\&R）出版，C 成为事实标准。
\end{timelinebox}

\chapter{K\&R C：事实标准的形成（1978-1989）}

\section{《The C Programming Language》}

\yrlbl{1978} Brian Kernighan（贝尔实验室的作家、UNIX 文档的主要作者之一）与 Dennis Ritchie 合著的 \textbf{《The C Programming Language}》由 Prentice-Hall 出版。这本书后来被简称为 \textbf{K\&R}（Kernighan \& Ritchie），成为 C 程序员的圣经。

K\&R 的影响力不仅在于教学，更在于它\textbf{事实上定义了 C 的语法和语义}。在 1989 年 ANSI C 标准发布之前的十年里，K\&R 第一版几乎就是 C 语言规范——任何 C 编译器都以 K\&R 的描述为基准，任何关于 C 行为的争议都以 K\&R 的文字为依据。

\begin{historybox}
\textbf{K\&R 的"Hello, World"}

K\&R 第一版第二章的第一个例子是：

\clabel\ \begin{lstlisting}[style=cstyle, numbers=none]
#include <stdio.h>

main()
{
    printf("hello, world\n");
}
\end{lstlisting}

注意几点：(1) \kw{main} 没有返回类型——在 K\&R C 中省略返回类型默认是 \kw{int}；(2) \kw{main} 没有参数列表——K\&R C 允许无原型函数声明；(3) \kw{printf} 不需要 \texttt{<stdio.h>} 的强制包含——只是约定。这三个"省略"在 ANSI C 中都被收紧：返回类型必须显式声明，函数必须有原型，\kw{printf} 必须包含 \texttt{<stdio.h>}。

"Hello, World" 本身并不是 K\&R 发明的——它来自 Kernighan 早年的 BCPL 教程。但通过 K\&R，它成为几乎所有编程教材的开篇约定。
\end{historybox}

\section{K\&R C 的语言特征}

K\&R C（即 1978-1989 年间事实上的 C）与现代 C（ANSI C / ISO C）有几处显著差异：

\begin{itemize}
  \item \textbf{函数原型缺失}：K\&R C 的函数声明只指定函数名，不指定参数类型。\kw{int f();} 表示"f 是一个返回 \kw{int} 的函数"，参数类型完全未指定——调用方传什么参数编译器都不报错。
  \item \textbf{默认 \kw{int}}：变量或函数没有显式类型时，默认为 \kw{int}。\kw{main()} 等价于 \kw{int main()}，\kw{f()} 等价于 \kw{int f()}。
  \item \textbf{\kw{register} 关键字}：提示编译器把变量放在 CPU 寄存器里。现代编译器普遍忽略这个提示。
  \item \textbf{无 \kw{const}}：K\&R 第一版没有 \kw{const} 关键字（虽然 Ritchie 在 1979 年的一个内部备忘录中提到过这个想法）。
  \item \textbf{无 \kw{void}}：早期 K\&R C 没有 \kw{void} 类型——函数要么返回 \kw{int}，要么不返回任何值（写法仍是 \kw{int f()}，但调用方忽略返回值）。\kw{void} 是在 K\&R 第一版后期和第二版中逐步加入的。
  \item \textbf{无 \kw{enum}}：枚举类型在 K\&R 第一版中没有，是后来扩展的。
\end{itemize}

\clabel\ \begin{lstlisting}[style=cstyle]
/* K&R 风格的函数定义（"old-style"） */
int max(a, b, c)
int a, b;       /* 参数类型在函数体外声明 */
char c;
{
    int m = a > b ? a : b;
    return m > c ? m : c;
}

/* ANSI C 风格的函数原型（C++ 引入，C89 采纳） */
int max(int a, int b, char c);
\end{lstlisting}

\section{UNIX 生态的扩张}

\yrlbl{1979} UNIX V7 发布，是第一个"可移植"的 UNIX 版本。Bell Labs 把 V7 源代码授权给大学与商业机构，催生了大量分支：

\begin{itemize}
  \item \textbf{BSD}（Berkeley Software Distribution，1977-1995）：加州大学伯克利分校开发，加入了 TCP/IP 协议栈、\kw{vi} 编辑器、csh shell。BSD 是后来 FreeBSD、NetBSD、OpenBSD、macOS（Darwin）的祖先。
  \item \textbf{System III / System V}（1982-）：AT\&T 商业化的 UNIX 分支，与 BSD 形成"两个 UNIX"的局面。
  \item \textbf{Xenix}（1980-1989）：Microsoft 授权的 UNIX 版本，运行在 Intel 8086/8088 上。是早期 PC 上最常见的 UNIX。
  \item \textbf{SunOS}（1982-）：Sun Microsystems 基于 BSD 开发，后来演变为 Solaris。
  \item \textbf{HP-UX}、\textbf{AIX}、\textbf{Ultrix}、\textbf{IRIX}：各硬件厂商的 UNIX 变体。
\end{itemize}

每一个 UNIX 变体都需要 C 编译器。Bell Labs 把 V7 的 C 编译器源代码也一并授权，于是各厂商纷纷基于这份源代码做自己的扩展——这是 C 语言"方言化"问题的开端。不同 UNIX 上的 C 在细节上各有差异：结构体对齐、整数大小、字节序、库函数名称都不完全一致。

\yrlbl{1983} 这种碎片化让"可移植 C 程序"成为一个理论概念而非实践。开发者抱怨"在 BSD 上能编译，在 System V 上就出错"。一些机构开始尝试写"portability guides"，但都不够权威。\textbf{C 语言需要一个正式标准}——这个共识在 1983 年前后形成。

\section{C 的"方言战争"与标准化的需求}

\yrlbl{1982-1985} 这一时期 C 的几个主要方言：

\begin{itemize}
  \item \textbf{K\&R C}：贝尔实验室的"原版"，最保守。
  \item \textbf{BSD C}：加入了 \kw{const}、\kw{volatile}、\kw{signed}、函数原型（受 C++ 影响）。
  \item \textbf{System V C}：与 BSD 类似但略有不同。
  \item \textbf{Microsoft C}（MS-C，1983-）：运行在 MS-DOS 与 Xenix 上，加入了大量 Microsoft 特有的扩展（\kw{\_far}、\kw{\_near}、\kw{\_pascal} 等内存模型修饰符）。
  \item \textbf{Borland Turbo C}（1987-）：加入了内联汇编等扩展。
\end{itemize}

方言之间的不兼容性已经严重到影响实际开发。一个跨平台 C 程序需要写大量 \kw{\#ifdef} 来处理差异。这正是 1983 年 ANSI 决定立项 C 标准的直接动因。

\begin{compare}
\textbf{K\&R C 与 ANSI C 的核心差异}

\kw{K\&R C}（1978-1989，事实标准）：函数声明可以无原型；返回类型默认 \kw{int}；无 \kw{const}、\kw{volatile}、\kw{signed} 关键字；无 \kw{void*}（用 \kw{char*} 替代）；无函数指针的标准转换规则；无标准库（只有约定俗成的函数名）。

\kw{ANSI C}（1989）：函数原型强制（来自 C++）；返回类型必须显式；新增 \kw{const}、\kw{volatile}、\kw{signed}、\kw{void*}；标准化所有库函数（\texttt{<stdio.h>}、\texttt{<stdlib.h>}、\texttt{<string.h>} 等 15 个头文件）；引入"严格符合"（strictly conforming）与"符合"（conforming）的区分。

差异的影响：K\&R C 程序大多能在 ANSI C 编译器上运行（编译器默认开启宽松模式），但 ANSI C 的\kw{const}、函数原型、\kw{void*} 等改进让程序更安全、更易维护。
\end{compare}

\chapter{ANSI X3J11 与 C89/C90：第一个国际标准}

\section{X3J11 委员会的成立}

\yrlbl{1983} 美国国家标准学会（ANSI，American National Standards Institute）成立了 \textbf{X3J11 委员会}，专门为 C 语言制定一个正式的美国国家标准。X3J11 的命名规则来自 ANSI 的项目编号系统——X3 是"信息技术"大类，J11 是 C 语言项目。

X3J11 的组成非常多元，体现了"利益相关方共同参与"的原则：

\begin{itemize}
  \item 编译器厂商：AT\&T、Microsoft、Borland、MetaWare、Mark Williams Company、Lattice 等。
  \item 用户代表：IBM、Apple、Sun、HP、DEC、Boeing、AT\&T Bell Labs。
  \item 学术界：多位大学教授。
  \item 独立顾问：包括 Kernighan 与 Ritchie 本人（虽然他们的角色是咨询而非投票）。
  \item 政府代表：NIST（美国国家标准与技术研究院）、DoD（美国国防部，他们要求 C 满足 MIL-STD 等政府标准）。
\end{itemize}

\yrlbl{1983-1988} X3J11 委员会花了五年多时间起草标准。这段时间委员会成员通过邮件列表、定期会议、提案投票来推进工作。这种"邮件 + 会议 + 投票"的工作模式后来被 ISO/IEC 委员会继承，成为现代编程语言标准化的标准范式。

\section{X3J11 的关键设计决策}

X3J11 在五年中做出的关键决策：

\begin{enumerate}
  \item \textbf{采纳 C++ 的函数原型}：函数声明必须指定参数类型，旧式 K\&R 函数定义虽然保留兼容，但新代码必须用原型。这是 C 标准化最重要的语义改进。
  \item \textbf{采纳 \kw{const} 与 \kw{volatile}}：从 C++ 借鉴的关键字，让程序员能表达"只读"与"硬件相关"的意图。
  \item \textbf{引入 \kw{void*} 作为通用指针}：K\&R C 用 \kw{char*} 充当通用指针，但类型不安全。ANSI C 引入 \kw{void*}，并定义所有对象指针都能隐式转换为 \kw{void*} 而不丢失信息。
  \item \textbf{标准化整个 C 运行时库}：15 个标准头文件（\texttt{<stdio.h>}、\texttt{<stdlib.h>}、\texttt{<string.h>}、\texttt{<math.h>} 等），约 150 个函数。每个函数的签名、行为、错误码都被严格定义。
  \item \textbf{定义"实现定义行为"（implementation-defined）}：标准明确划分了哪些行为是"标准要求的"、哪些是"实现定义的"（编译器必须文档化）、哪些是"未指定的"（unspecified）、哪些是"未定义行为"（UB）。这套分类后来成为所有 C/C++ 标准文档的基础。
  \item \textbf{拒绝面向对象扩展}：当时 C++ 已经存在，但 X3J11 明确拒绝把 class、继承等概念纳入 C 标准。这个决策让 C 保持了"小语言"的定位，与 C++ 形成清晰的分工。
  \item \textbf{保留向后兼容性}：所有合法的 K\&R C 程序都应该是合法的 ANSI C 程序（除了一些边缘情况）。这个原则贯穿后续所有 C 标准。
\end{enumerate}

\section{C89 与 C90：从 ANSI 到 ISO}

\yrlbl{1989} X3J11 完成标准草案，ANSI 正式发布 \textbf{ANSI X3.159-1989}，俗称 \textbf{C89} 或 \textbf{ANSI C}。这是 C 的第一个正式标准。

\yrlbl{1990} ISO/IEC JTC1/SC22 几乎原封不动地采纳了 ANSI X3.159-1989，发布为 \textbf{ISO/IEC 9899:1990}，俗称 \textbf{C90}。两份标准在技术内容上完全相同，仅在编辑格式与少量措辞上有差异。从此"C89"与"C90"成为同一标准的两个名称——美国称 C89，国际称 C90。

\begin{stdbox}
\textbf{ISO/IEC 9899:1990 的标准结构}

标准分为若干条款（clause），每个条款编号是后续 C 标准都遵循的：

\kw{1} 范围（Scope）；\kw{2} 术语定义（Definitions）；\kw{3} 概念（Concepts）；\kw{4} 符合性（Conformance）；\kw{5} 环境（Environment）；\kw{6} 语言（Language）；\kw{7} 标准库（Standard Library）；\kw{8} 移植性考虑（Portability）。

其中条款 \kw{6} 是语言核心规范，按词法、语法、语义、约束子条款组织。这种"语法 + 约束 + 语义 + 示例"的四段式结构沿用至今。
\end{stdbox}

\yrlbl{1994-1996} ISO 发布了 \textbf{Amendment 1}（ISO/IEC 9899 AMD1:1995），加入了一些国际化支持（多字节字符、宽字符库 \texttt{<wchar.h>}、\texttt{<wctype.h>}）和小幅修订。这份修订与 C90 合并后通常称为 \textbf{C95} 或 \textbf{C90/Amd1}。

\section{从 C90 到 C99：现代化}

\yrlbl{1999} 在 C90 发布将近十年之后，ISO 终于发布了重大更新版本 \textbf{ISO/IEC 9899:1999}，俗称 \textbf{C99}。C99 的更新内容反映了一个十年间编译器实践与用户需求的积累：

\begin{itemize}
  \item \textbf{\kw{//} 单行注释}：从 C++ 借来，C90 不允许。
  \item \textbf{变量声明位置自由化}：C90 要求所有变量必须在块的开头声明，C99 允许在块的任意位置声明。
  \item \textbf{\kw{inline} 关键字}：标准化了各编译器都有的内联提示。
  \item \textbf{变长数组（VLA，Variable-Length Array）}：\kw{int a[n];} 其中 \kw{n} 是运行期变量。这是 C99 最具争议的特性，C11 把它降级为"可选"。
  \item \textbf{复合字面量（compound literals）}：\kw{(int[])\{1,2,3\}}，允许在表达式中创建匿名数组。
  \item \textbf{指定初始化器（designated initializers）}：\kw{struct Point p = \{.x = 1, .y = 2\};}。
  \item \textbf{\kw{long long} 类型}：64 位整数，至少 64 位宽。
  \item \textbf{\kw{\_Bool} 类型}：标准布尔类型（在 \texttt{<stdbool.h>} 中提供 \kw{bool} 宏）。
  \item \textbf{\kw{\_Complex} 与 \kw{\_Imaginary}}：复数与虚数支持（在 \texttt{<complex.h>} 中）。
  \item \textbf{\kw{restrict} 指针限定符}：告诉编译器该指针是访问其指向对象的唯一方式，便于优化。
  \item \textbf{\kw{<stdint.h>} 与 \kw{<inttypes.h>}}：精确宽度的整数类型（\kw{int32\_t}、\kw{uint64\_t} 等）与格式化宏。
  \item \textbf{\kw{<tgmath.h>}}：泛型数学宏，根据参数类型自动选择 \kw{sin}、\kw{sinf}、\kw{sinl}。
  \item \textbf{灵活的数组成员（FAM，Flexible Array Member）}：\kw{struct S \{ int n; int data[]; \};} 允许结构体最后一个成员是"零长数组"，与后续分配配合实现变长记录。
  \item \textbf{\kw{<wchar.h>} 增强}：宽字符 I/O、转换函数。
\end{itemize}

\clabel\ \begin{lstlisting}[style=cstyle]
/* C99 新特性示例 */
#include <stdio.h>
#include <stdint.h>
#include <stdbool.h>

// 单行注释（C99 新增）
struct Point { int x, y; };

int main(void) {
    // 变量声明在任意位置（C99 新增）
    int n = 5;

    // VLA（C99 新增）
    int dynamic_array[n];

    // 复合字面量（C99 新增）
    struct Point p1 = (struct Point){ .x = 1, .y = 2 };

    // C99 整数类型
    int32_t i32 = INT32_MAX;
    uint64_t u64 = UINT64_C(18446744073709551615);

    // C99 bool
    bool flag = true;

    printf("%d %jd %ju %d\n", p1.x, (intmax_t)i32, (uintmax_t)u64, flag);
    return 0;
}
\end{lstlisting}

\section{C99 的采用困境}

C99 标准虽然技术上很完整，但实际采用却比预期慢得多。这反映了\textbf{去中心化标准化的一个根本特点}：标准本身只是文档，要让它真正发挥作用，必须依赖编译器实现与库支持。

\begin{itemize}
  \item \textbf{GCC}：到 GCC 4.5（2010）才完整支持 C99，包括 VLA 与 FAM。
  \item \textbf{Clang}：从一开始就完整支持 C99（除少量边缘特性）。
  \item \textbf{Microsoft Visual C++}（MSVC）：长期不支持 C99。MSVC 2013 才开始支持 \texttt{<stdint.h>}，MSVC 2015 才支持 \kw{\_Bool}，VLA 至今不支持。
  \item \textbf{嵌入式编译器}（IAR、Keil、ARM Compiler）：支持程度参差不齐，许多旧版本只支持 C90。
\end{itemize}

MSVC 的 C99 支持滞后影响最大。Microsoft 长期把精力投入 C++，C 编译器只是 C++ 编译器的子集——MSVC 的 \kw{cl.exe} 编译 \kw{.c} 文件时使用一个非常老的 C 前端。这种"编译器厂商选择性实现"的现象，正是后面要讨论的\textbf{去中心化治理缺点}的一个具体体现：ISO 发布了标准，但没有任何机制强迫编译器实现。

\chapter{C11 → C23：现代化的延续}

\section{C11：并发与原子}

\yrlbl{2011} ISO/IEC 9899:2011，俗称 \textbf{C11}（最初命名为 C1X，X 是占位符）。C11 的最大特性是把并发与原子操作引入标准 C：

\begin{itemize}
  \item \textbf{\kw{\_Atomic} 类型限定符}：原子操作语义，标准头文件 \seqkw{<stdatomic.h>} 提供完整接口；与 C++11 的 \seqkw{std::atomic} 共享同一个内存模型。
  \item \textbf{\kw{\_Thread\_local} 存储类}：线程局部存储（TLS），\texttt{<threads.h>} 提供完整线程库（\kw{thrd\_create}、\kw{mtx\_lock}、\kw{cnd\_wait} 等）。
  \item \textbf{\kw{\_Generic} 选择器}：编译期类型分发，类似 C++ 的函数重载。
  \item \textbf{\kw{\_Static\_assert} 与 \kw{\_Alignas} / \kw{\_Alignof}}：从 C++ 借来的编译期断言与对齐控制。
  \item \textbf{\kw{\_Noreturn} 函数限定符}：标记永不返回的函数（如 \kw{exit}）。
  \item \textbf{边界检查接口}：Annex K 定义了一组 \seqkw{\_s} 后缀的安全版函数，如 \seqkw{strcpy\_s} 与 \seqkw{sprintf\_s}，但实现极少，C23 已建议废弃。
  \item \textbf{匿名结构体/联合成员}：\kw{struct \{ union \{ int i; float f; \}; \};} 允许内部 union 没有名字。
  \item \textbf{VLA 降级为可选}：C99 把 VLA 列为强制特性，但实际编译器（特别是 MSVC）支持差，C11 把 VLA 改为"可选实现"。
  \item \textbf{\kw{<uchar.h>}}：UTF-16（\kw{char16\_t}）与 UTF-32（\kw{char32\_t}）支持。
\end{itemize}

\clabel\ \begin{lstlisting}[style=cstyle]
/* C11 示例：原子操作与泛型选择器 */
#include <stdatomic.h>
#include <stdio.h>
#include <threads.h>

_Atomic int counter = 0;

int increment(void* arg) {
    (void)arg;
    for (int i = 0; i < 1000; ++i) {
        atomic_fetch_add(&counter, 1);  /* 原子递增 */
    }
    return 0;
}

/* _Generic：编译期类型分发 */
#define print_value(x) _Generic((x), \
    int:    printf("%d\n", x),       \
    double: printf("%f\n", x),       \
    char*:  printf("%s\n", x),       \
    default: printf("unknown\n"))

int main(void) {
    thrd_t t1, t2;
    thrd_create(&t1, increment, NULL);
    thrd_create(&t2, increment, NULL);
    thrd_join(t1, NULL);
    thrd_join(t2, NULL);

    print_value(42);          /* 走 int 分支 */
    print_value(3.14);        /* 走 double 分支 */
    print_value("hello");     /* 走 char* 分支 */

    printf("counter = %d\n", atomic_load(&counter));
    return 0;
}
\end{lstlisting}

\section{C17：纯修订版}

\yrlbl{2018} ISO/IEC 9899:2018，俗称 \textbf{C17}（也叫 C18）。C17 没有引入任何新特性——它只是 C11 的\textbf{缺陷修订版}（defect report fixes）。这种"零特性修订"在 C 历史上是第一次，体现了委员会对稳定性的重视。

\section{C23：与 C++ 的同步}

\yrlbl{2024} ISO/IEC 9899:2024，俗称 \textbf{C23}（也叫 C2X）。C23 是 C11 之后最重要的更新，引入了一系列与 C++ 同步的特性：

\begin{itemize}
  \item \textbf{\kw{bool}、\kw{true}、\kw{false} 成为关键字}：C99 通过 \texttt{<stdbool.h>} 提供宏，C23 直接成为语言内置（与 C++23 同步）。
  \item \textbf{\kw{nullptr} 关键字}：与 C++11 的 \kw{nullptr} 一致，是类型安全的空指针常量。
  \item \textbf{\kw{typeof} 与 \kw{typeof\_unqual}}：类型推导（GCC 长期作为扩展提供，C23 标准化）。
  \item \textbf{\kw{constexpr}}：编译期常量表达式（与 C++ 同名同义）。
  \item \textbf{属性语法 \kw{[[...]]}}：与 C++11 的 \kw{[[...]]} 一致，C23 采纳 \kw{[[deprecated]]}、\kw{[[fallthrough]]}、\kw{[[maybe\_unused]]}、\kw{[[nodiscard]]}、\kw{[[noreturn]]} 等。
  \item \textbf{\kw{auto} 关键字语义变化}：C 早期的 \kw{auto} 表示"自动存储期"（即局部变量，默认就是），C23 把 \kw{auto} 重定义为类型推导（与 C++ 一致），原"自动存储期"语义被取消。
  \item \textbf{二进制字面量 \kw{0b1010}}：与 C++14 同步。
  \item \textbf{\kw{\_Decimal32} / \kw{\_Decimal64} / \kw{\_Decimal128}}：十进制浮点（IEC 60559 标准），用于金融计算。
  \item \textbf{\kw{<stdbit.h>}}：位操作工具（\kw{stdc\_count\_ones}、\kw{stdc\_leading\_zeros} 等）。
  \item \textbf{\kw{<stdckdint.h>}}：溢出检查的算术运算（\kw{ckd\_add}、\kw{ckd\_sub}、\kw{ckd\_mul}）。
  \item \textbf{函数定义中可省略 \kw{void} 参数}：\kw{int main() \{\}} 现在等价于 \kw{int main(void) \{\}}（与 C++ 一致）。
  \item \textbf{移除若干废弃特性}：\kw{gets()} 函数（早已被认为不安全）、\kw{<tgmath.h>} 部分宏、Annex K 的"\_s"接口被标记为废弃。
  \item \textbf{预定义宏 \seqkw{\_\_STDC\_VERSION\_\_} = \seqkw{202311L}}。
\end{itemize}

\clabel\ \begin{lstlisting}[style=cstyle]
/* C23 示例：与 C++ 同步的关键字 */
#include <stdio.h>
#include <stdlib.h>

int main() {                     /* 等价于 int main(void) */
    bool flag = true;            /* bool 是关键字（C23） */
    int* p = nullptr;            /* nullptr 是关键字（C23） */

    /* auto 类型推导（C23） */
    auto x = 3.14;               /* x 是 double */
    auto y = 42;                 /* y 是 int */

    /* typeof 类型推导（C23） */
    typeof(x) z = x * 2;         /* z 是 double */

    /* constexpr 编译期常量（C23） */
    constexpr int N = 10;

    /* 属性语法（C23） */
    [[maybe_unused]] int unused_var = 0;
    [[deprecated("use new_func instead")]] void old_func(void);

    /* 二进制字面量（C23） */
    int mask = 0b1010;

    printf("%d %g %d %d\n", flag, x, y, mask);
    return 0;
}
\end{lstlisting}

\begin{compare}
\textbf{C 与 C++ 的同步：C23 案例}

C23 是历史上第一次 C 标准明显向 C++ "看齐"：\kw{bool}、\kw{nullptr}、\kw{constexpr}、\kw{[[...]]} 属性、\kw{auto} 类型推导、\kw{typeof} 都直接来自 C++。

但同步并非完全——C23 没有引入 C++ 的类、模板、引用、运算符重载、命名空间、异常等核心特性。C 委员会（WG14）与 C++ 委员会（WG21）保持着"协作但独立"的关系：WG21 的某些提案（如 \kw{<bit>}、\kw{<source\_location>}）会先在 WG14 讨论是否有 C 版本；反之 C23 的 \kw{typeof} 也影响了 C++23 的 \kw{typeof}（虽然 C++ 已经有 \kw{decltype}）。

这种"两个独立委员会、同一组织、定期沟通"的模式是 ISO/IEC SC22 治理结构的特征之一，后面会详细讨论。
\end{compare}

% ====================================================================
\part{C++ 的诞生与标准化}
% ====================================================================

\chapter{C with Classes：1979-1983}

\section{Stroustrup 的 Simula 经验}

\yrlbl{1975} Bjarne Stroustrup 是丹麦人，在剑桥大学（University of Cambridge）取得博士学位。他的博士研究是分布式系统的模拟，使用了挪威开发的 \textbf{Simula 67}——第一门面向对象语言。Simula 的 \kw{class} 与继承机制让 Stroustrup 印象深刻，能在语言层面清晰表达"模拟实体之间的关系"。

但 Simula 的实现太慢——Simula 编译生成的代码效率不足以处理大规模模拟。Stroustrup 的博士论文不得不退回用 BCPL/C 来写关键部分，丢失了 Simula 的语义优雅。

\yrlbl{1979} Stroustrup 回到 Bell Labs，加入了 Thompson 与 Ritchie 的 UNIX 团队。他面临一个具体问题：UNIX 内核的分布式模拟需要同时具备 Simula 的类组织能力和 C 的执行效率。于是他启动了"\textbf{C with Classes}"项目——在 C 的基础上加入 Simula 风格的类机制。

\begin{historybox}
\textbf{为什么选 C 作为基础}

Stroustrup 后来解释过这个决策的几条原因：

(1) \textbf{性能}：C 是为系统编程设计的，没有运行期开销。

(2) \textbf{可移植性}：C 已经被移植到十几种机器上，C with Classes 也能跟着移植。

(3) \textbf{灵活性}：C 信任程序员，允许做"危险"的操作（指针算术、类型转换），这正适合系统编程。

(4) \textbf{与 UNIX 共生}：C 是 UNIX 的语言，Bell Labs 的所有系统软件都是 C 写的，C with Classes 可以直接集成。

(5) \textbf{没有更好的选择}：当时 BCPL 已经过时，Pascal 不适合系统编程（没有指针算术、没有底层访问），Algol 68 太复杂，Ada 还没出现。

这个决策的代价是：C with Classes 必须保留 C 的全部"危险特性"，包括无边界检查的数组、隐式类型转换、宏等。后来 C++ 一直被批评为"过于复杂"，根源就在这里。
\end{historybox}

\section{Cfront 与第一批特性}

\yrlbl{1979-1983} Stroustrup 实现了一个名为 \textbf{Cfront}（"C front"，意为"C 的前端"）的预处理器，把 C with Classes 的源码\textbf{翻译}成普通 C 代码，再用 C 编译器编译。这种"翻译到 C"的策略让 Cfront 能快速移植到任何有 C 编译器的平台。

C with Classes 的初版特性：

\begin{itemize}
  \item \textbf{\kw{class} 关键字}：封装数据与函数。
  \item \textbf{构造函数与析构函数}：自动初始化与清理。
  \item \textbf{继承（public/private）}：单继承，支持访问控制。
  \item \textbf{成员函数与访问修饰符}：\kw{public}、\kw{private}（\kw{protected} 后来加入）。
  \item \textbf{强类型检查}：相比 K\&R C，函数声明强制原型。
  \item \textbf{函数原型}：直接借用了 Stroustrup 之前在 Simula 中学到的概念，并影响 ANSI C 委员会采纳。
\end{itemize}

\cpplabel\ \begin{lstlisting}
/* C with Classes 风格代码（1980 年前后） */
class Stack {
    int data[100];
    int top;
public:
    Stack() { top = 0; }
    void push(int x) { data[top++] = x; }
    int pop() { return data[--top]; }
};

class BoundedStack : Stack {     /* 继承（早期语法，无冒号） */
    int max_size;
public:
    void push(int x) {
        if (top < max_size) Stack::push(x);
    }
};
\end{lstlisting}

\section{命名与第一批用户}

\yrlbl{1983} C with Classes 在 Bell Labs 内部已有约 50 名用户，分布在操作系统、网络协议、数据库等多个项目。Stroustrup 决定给这门语言一个更短、更"市场友好"的名字。候选包括"New C"、"C+++"、"C83"等，最终选定 \textbf{C++}——意为"C 的下一个版本"，与 C 的命名传统一致（B → C 也是按字母顺序）。

\yrlbl{1983} 同年，C++ 加入了若干关键特性：

\begin{itemize}
  \item \textbf{虚函数（virtual functions）}：支持运行期多态，这是面向对象的核心机制。
  \item \textbf{运算符重载}：\kw{operator+}、\kw{operator=} 等，让自定义类型能像内置类型一样使用。
  \item \textbf{引用（references）}：\kw{int\& r = x;}，避免显式指针的语法噪音。
  \item \textbf{\kw{const}}：与 C 的 \kw{const} 兼容，但语义更严格（C++ 的 \kw{const} 默认有内部链接性，C 默认有外部链接性）。
\end{itemize}

虚函数的加入标志着 C++ 从"带类的 C"真正变成了"面向对象语言"。在此之前，C++ 的多态只能通过函数指针手动实现。

\chapter{CFront 与 C++ 的商业化：1983-1989}

\section{AT\&T 发布 CFront 1.0}

\yrlbl{1985} AT\&T（Bell Labs 的母公司）正式向公众发布 \textbf{CFront 1.0}——C++ 的第一个商业编译器。同时 Stroustrup 出版了 \textbf{《The C++ Programming Language}》第一版（俗称 TC++PL）。这两件事标志着 C++ 从 Bell Labs 内部工具变成了公开语言。

CFront 1.0 的实现策略：

\begin{itemize}
  \item 把 C++ 源码翻译为 C 源码，然后用任何符合 K\&R C 的编译器编译。
  \item 翻译过程包括：类 → 结构体 + 函数；虚函数 → 虚函数表（vtable）；运算符重载 → 命名函数。
  \item 这种策略让 CFront 能快速移植——只需要目标平台有 C 编译器，CFront 就能工作。
\end{itemize}

CFront 的源码以"非独家许可"的方式授权给数十家公司。1985-1990 年间，市场上出现了大量基于 CFront 的 C++ 编译器：Microsoft C++ 1.0、Borland Turbo C++、Zortech C++、Symantec C++ 等。

\section{TC++PL 与 ARM}

\yrlbl{1985} TC++PL 第一版出版，迅速成为 C++ 程序员的标准参考。这本书与 K\&R 类似，事实上定义了 C++ 的语言规范——在 1998 年 ISO C++ 发布之前，TC++PL 是 C++ 行为的权威说明。

\yrlbl{1990} Margaret Ellis 与 Stroustrup 合著的 \textbf{《The Annotated C++ Reference Manual}》（俗称 \textbf{ARM}，Annotated Reference Manual）出版。ARM 详细描述了 C++ 2.1 版本的语法、语义、约束，并为每个规则附上"为什么这样设计"的注释。ARM 后来成为 ANSI/ISO C++ 委员会的起点参考文档。

\section{1989 年的关键提案：模板与异常}

\yrlbl{1989} Stroustrup 在 USENIX C++ 大会上提交了一份里程碑式的提案 \textbf{《Parameterized Types for C++}》，提议给 C++ 加入\textbf{模板}（templates）——一种编译期泛型机制。这份提案引发了长达两年的讨论：

\begin{itemize}
  \item 模板的核心想法是\textbf{编译期代码生成}：写一份代码，编译器为不同类型生成多份。
  \item 反对意见：模板会大幅膨胀编译产物（"code bloat"），让二进制变大、编译变慢。
  \item 支持意见：没有模板就无法实现类型安全的容器。考虑 \seqkw{vector<int>}：它和 \seqkw{vector<double>} 必须各自生成独立代码，无法共用一份实现。
\end{itemize}

\yrlbl{1990} 模板被 C++ 2.0 接受。同年加入的还有\textbf{异常处理}（\kw{try}、\kw{catch}、\kw{throw}）与\textbf{运行期类型识别}（RTTI，\kw{dynamic\_cast}、\kw{typeid}）。这三个特性——模板、异常、RTTI——共同构成了 C++ "现代风格"的基础。

C++ 2.0 的其他关键特性：

\begin{itemize}
  \item \textbf{多重继承}（multiple inheritance）：允许一个类继承自多个基类。
  \item \textbf{抽象类}（abstract classes）与\textbf{纯虚函数}（\kw{virtual ... = 0}）。
  \item \textbf{\kw{protected} 访问修饰符}。
  \item \textbf{成员指针}（pointer to member）。
  \item \textbf{命名空间}（namespace）的早期形式。
\end{itemize}

\cpplabel\ \begin{lstlisting}
/* C++ 2.0 风格（1989-1991）：模板 + 异常 + 多重继承 */
template <class T>           /* 早期模板语法用 class，不用 typename */
class Vector {
    T* data;
    int size;
public:
    Vector(int n) : data(new T[n]), size(n) {}
    ~Vector() { delete[] data; }
    T& operator[](int i) {
        if (i < 0 || i >= size) throw "out of range";   /* 异常 */
        return data[i];
    }
};

class Shape {
public:
    virtual void draw() = 0;        /* 纯虚函数 */
};

class Drawable {
public:
    virtual bool canDraw() { return true; }
};

class Circle : public Shape, public Drawable {  /* 多重继承 */
public:
    void draw() override { /* ... */ }
};
\end{lstlisting}

\section{编译器市场的爆发}

\yrlbl{1985-1991} C++ 商业化期间，市场上涌现了多家 C++ 编译器：

\begin{itemize}
  \item \textbf{Microsoft C/C++}（1985-）：MS-DOS、Windows 平台的主流编译器。
  \item \textbf{Borland Turbo C++}（1990-）：以快速编译、集成 IDE、低价著称，对 PC 编程教学影响巨大。
  \item \textbf{Zortech C++}（1988-）：第一家发布"原生 C++ 编译器"（不通过 CFront 翻译到 C）的厂商。
  \item \textbf{Symantec C++}（1989-）：收购 Zortech 后改名。
  \item \textbf{Sun C++}（1988-）：SPARC/Solaris 平台。
  \item \textbf{HP aC++}（1990-）：HP-UX 平台。
  \item \textbf{IBM C++/System}（1991-）：AIX、OS/2 平台。
\end{itemize}

不同编译器之间的差异比 C 时代更严重——C++ 没有事实标准，每家公司都按自己对 ARM 的理解实现。同一个 C++ 程序在 Microsoft、Borland、Sun 上行为不同是常态。这种碎片化是 ANSI/ISO 立项 C++ 标准化的直接动因。

\chapter{C++98 / C++03：第一次国际标准}

\section{ANSI/ISO C++ 委员会成立}

\yrlbl{1991} ANSI 成立了 \textbf{X3J16 委员会}（C 是 X3J11，C++ 是 X3J16），负责 C++ 标准化。

\yrlbl{1993} ISO/IEC JTC1/SC22 成立了 \textbf{WG21}（Working Group 21），专门负责 C++ 国际标准。WG21 与 X3J16 合并工作——同一组人、同一个会议、同一份邮件列表，但产出两份标准（ANSI 与 ISO）。这种"ANSI/ISO 联合委员会"模式后来被许多编程语言效仿。

WG21 的工作量极大：

\begin{itemize}
  \item 需要把 ARM 中模糊的语言定义转为精确的标准措辞。
  \item 需要协调编译器厂商之间已经存在的实现差异。
  \item 需要设计一个\textbf{标准库}——C++ 之前没有官方标准库，只有 AT\&T 的 \kw{libg++}、各厂商各自的库。
\end{itemize}

\section{STL 的纳入}

\yrlbl{1994} Alexander Stepanov（HP 实验室）与 Meng Lee 提出了 \textbf{STL}（Standard Template Library，标准模板库）的设计。STL 的核心思想：

\begin{itemize}
  \item \textbf{容器}（containers）：\kw{vector}、\kw{list}、\kw{deque}、\kw{map}、\kw{set} 等。
  \item \textbf{迭代器}（iterators）：统一的访问接口，让算法不依赖具体容器。
  \item \textbf{算法}（algorithms）：\kw{sort}、\kw{find}、\kw{copy}、\kw{transform} 等通用算法。
  \item \textbf{函数对象}（functors）与\textbf{适配器}（adaptors）：可组合的函数式编程工具。
  \item \textbf{分配器}（allocators）：抽象内存分配策略。
\end{itemize}

STL 是 C++ 模板编程的集大成者——它证明了"通用容器 + 通用算法 + 通用迭代器"可以用模板优雅实现，没有性能损失。

\yrlbl{1994} WG21 投票决定把 STL 纳入 C++ 标准库。这是 C++ 标准化过程中最大的单一决策——它让标准库从一个"基础工具集合"变成了一个"通用编程范式"。

\section{C++98 发布}

\yrlbl{1998} ISO/IEC 14882:1998 发布，俗称 \textbf{C++98}。这是 C++ 的第一个国际标准，也是当时\textbf{最大的编程语言标准}——800 多页，远超 C90（约 200 页）。

C++98 的关键内容：

\begin{itemize}
  \item \textbf{语言核心}：基于 ARM，但加入了若干新特性与澄清：
    \begin{itemize}
      \item \kw{typename} 关键字（用于模板中区分类型与变量）。
      \item \kw{bool} 类型（独立于 \kw{int}）。
      \item \kw{explicit} 构造函数（防止隐式转换）。
      \item \kw{mutable} 成员（在 \kw{const} 方法中可修改）。
      \item \kw{static\_cast}、\kw{dynamic\_cast}、\kw{const\_cast}、\kw{reinterpret\_cast}（四种命名的转换运算符，替代 C 风格的强制转换）。
      \item \kw{true} 与 \kw{false} 字面量。
      \item 模板特化（template specialization）。
      \item 命名空间（namespace）正式标准化。
    \end{itemize}
  \item \textbf{标准库}：以 STL 为核心，加入了 I/O 流（\texttt{<iostream>}）、字符串（\texttt{<string>}）、数值（\texttt{<complex>}、\texttt{<valarray>}）、本地化（\texttt{<locale>}）、异常（\texttt{<exception>}、\texttt{<stdexcept>}）、诊断（\texttt{<cassert>}）、C 运行时兼容（\texttt{<cstdio>} 等）。
  \item \textbf{没有线程库}：C++98 标准库不包含线程支持，依赖 POSIX 或平台特定 API。这个缺失直到 C++11 才补上。
\end{itemize}

\section{C++03：技术勘误版}

\yrlbl{2003} ISO/IEC 14882:2003，俗称 \textbf{C++03}。与 C17 类似，C++03 是 C++98 的\textbf{缺陷修订版}（Technical Corrigendum 1），没有引入新特性。主要修复了 C++98 标准文本中约 200 处歧义、错误描述与不一致。

C++98/C++03 通常合称为 \textbf{C++98/03}，代表"C++11 之前的传统 C++"。

\section{委员会休眠期与 C++0x 项目}

\yrlbl{2003-2011} C++03 发布后，WG21 进入了一段相对缓慢的时期。委员会开始规划下一个重大版本——\textbf{C++0x}（x 是占位符，原本预期 200X 年发布，最终拖到 2011）。

C++0x 项目期间，委员会做了大量准备工作：

\begin{itemize}
  \item 收集编译器实现中的常见问题（Defect Reports）。
  \item 讨论语言演化方向（向后兼容 vs 新特性）。
  \item 评估各类提案的影响范围。
  \item 与编译器厂商协调实验性实现。
\end{itemize}

这段"休眠期"也暴露了去中心化标准化的一个特点：\textbf{没有外部压力时，委员会容易陷入冗长讨论}。C++0x 期间有大量提案被反复讨论数年，最终因共识不足被否决（如 Concepts 第一版、Garbage Collection 支持、Modules 早期版本）。这种"慢"既是优点（避免草率决策）也是缺点（无法快速响应社区需求）。

\chapter{现代 C++：从 C++11 到 C++26}

\section{C++11：一门新语言}

\yrlbl{2011} ISO/IEC 14882:2011 发布，俗称 \textbf{C++11}（也叫 C++0x）。这是 C++ 历史上最大的一次更新，Stroustrup 称之为"\textbf{一门新语言}"——如此多的新特性让 C++11 与 C++98 看起来像两种语言。

C++11 的核心特性按类别归纳：

\begin{itemize}
  \item \textbf{移动语义与右值引用}：\kw{int\&\& r} 与 \kw{std::move}，解决长期困扰 C++ 的"复制开销"问题。让 \kw{std::vector} 等容器在重新分配时能"移动"而非"复制"。
  \item \textbf{\kw{auto} 类型推导}：\kw{auto x = expr;} 让编译器推导类型，减少冗长声明。
  \item \textbf{\kw{decltype}}：获取表达式的类型，用于模板元编程。
  \item \textbf{Lambda 表达式}：\kw{[](int x) \{ return x * 2; \}}，闭包函数。
  \item \textbf{\kw{nullptr}}：类型安全的空指针，替代 \kw{NULL} 与字面量 \kw{0}。
  \item \textbf{智能指针}：\kw{std::unique\_ptr}、\kw{std::shared\_ptr}、\kw{std::weak\_ptr}，引入 RAII 风格的内存管理。
  \item \textbf{\kw{constexpr}}：编译期函数与对象，让 C++ 能做编译期计算。
  \item \textbf{\kw{enum class}}：强类型枚举，避免传统 \kw{enum} 的隐式转换问题。
  \item \textbf{范围 \kw{for}}：\kw{for (auto\& x : container)}。
  \item \textbf{变参模板}：\kw{template<typename... Args>}，支持任意数量模板参数。
  \item \textbf{初始化列表}：\kw{std::vector<int> v = \{1, 2, 3\};}，\kw{std::initializer\_list}。
  \item \textbf{委托构造、继承构造、\kw{override} / \kw{final}}。
  \item \textbf{\kw{noexcept}}：异常规范，让编译器优化无异常路径。
  \item \textbf{\kw{static\_assert}}：编译期断言。
  \item \textbf{\kw{thread\_local}}：线程局部存储。
  \item \textbf{标准库新增}：\texttt{<thread>}、\texttt{<mutex>}、\texttt{<atomic>}、\texttt{<future>}、\texttt{<condition\_variable>}（首次有线程支持）、\texttt{<chrono>}（时间库）、\texttt{<tuple>}、\texttt{<array>}、\texttt{<unordered\_map>}、\texttt{<unordered\_set>}、\texttt{<random>}、\texttt{<regex>}、\texttt{<system\_error>}。
\end{itemize}

\cppxxlabel{C++11}\ \begin{lstlisting}
/* C++11 风格代码示例 */
#include <vector>
#include <memory>
#include <thread>
#include <mutex>
#include <chrono>
#include <algorithm>

std::mutex m;
std::vector<int> shared_data;

void worker(int id) {
    auto data = std::make_unique<int>(id * 10);   // 智能指针 + auto
    std::lock_guard<std::mutex> lock(m);          // RAII 锁
    shared_data.push_back(*data);                 // 移动语义自动启用
}

int main() {
    std::vector<std::thread> threads;
    for (int i = 0; i < 4; ++i) {
        threads.emplace_back(worker, i);          // emplace + lambda 风格
    }
    for (auto& t : threads) t.join();             // 范围 for

    // Lambda + STL 算法
    std::sort(shared_data.begin(), shared_data.end(),
              [](int a, int b) { return a > b; });

    // constexpr 编译期计算
    constexpr int squared = [](int x) constexpr { return x * x; }(5);
    static_assert(squared == 25, "compile-time check");

    // chrono 时间库
    auto start = std::chrono::steady_clock::now();
    auto elapsed = std::chrono::steady_clock::now() - start;
    return 0;
}
\end{lstlisting}

\section{C++14：小修订}

\yrlbl{2014} ISO/IEC 14882:2014，俗称 \textbf{C++14}。C++14 是 C++11 的小修订版，主要新增：

\begin{itemize}
  \item \textbf{泛型 lambda}：\kw{[](auto x) \{ return x * 2; \}}，lambda 参数可用 \kw{auto}。
  \item \textbf{函数返回类型推导}：\kw{auto f() \{ return 42; \}}，函数返回类型可用 \kw{auto} 或 \kw{decltype(auto)}。
  \item \textbf{变量模板}：\kw{template<typename T> constexpr T pi = T(3.14159);}。
  \item \textbf{放松的 \kw{constexpr}}：C++11 的 \kw{constexpr} 函数只能有一条返回语句，C++14 允许循环、分支、局部变量。
  \item \textbf{二进制字面量}：\kw{0b1010}。
  \item \textbf{数字分隔符}：\kw{1'000'000}。
  \item \textbf{\kw{std::make\_unique}}：补齐 C++11 的 \kw{make\_shared} 缺失。
  \item \textbf{\kw{[[deprecated]]} 属性}：标记废弃接口。
\end{itemize}

\section{C++17：实用主义更新}

\yrlbl{2017} ISO/IEC 14882:2017，俗称 \textbf{C++17}。C++17 是"实用主义"的更新，引入了大量日常编程便利特性：

\begin{itemize}
  \item \textbf{结构化绑定}：\kw{auto [x, y, z] = tuple;}，解构 tuple、pair、struct、array。
  \item \textbf{\kw{if constexpr}}：编译期分支，让模板代码更清晰。
  \item \textbf{\kw{if}/\kw{switch} 初始化语句}：\kw{if (auto it = m.find(k); it != m.end()) \{...\}}。
  \item \textbf{折叠表达式}：\kw{(args + ...)}，简化变参模板。
  \item \textbf{类模板参数推导}（CTAD）：\kw{std::pair p(1, 2.0);} 不需要写 \kw{<int, double>}。
  \item \textbf{内联变量}：\kw{inline int x = 0;} 可以在头文件中定义而不引起多重定义错误。
  \item \textbf{嵌套命名空间}：\kw{namespace A::B::C \{\}}。
  \item \textbf{\kw{[[nodiscard]]}、\kw{[[maybe\_unused]]}、\kw{[[fallthrough]]} 属性}。
  \item \textbf{标准库新增}：
    \begin{itemize}
      \item \kw{std::optional}、\kw{std::variant}、\kw{std::any}（"vocabulary types"）。
      \item \kw{std::string\_view}（非拥有的字符串视图）。
      \item \kw{std::filesystem}（文件系统库，从 Boost.Filesystem 移植）。
      \item \kw{std::invoke}、\kw{std::apply}。
      \item 并行算法：\kw{std::execution::par}、\kw{std::execution::seq}（执行策略）。
    \end{itemize}
  \item \textbf{移除}：\kw{std::auto\_ptr}（被 \kw{unique\_ptr} 替代）、\kw{std::random\_shuffle}、\kw{throw()} 动态异常规范（被 \kw{noexcept} 替代）、三元运算符注册字符（\kw{??=} 等）。
\end{itemize}

\cppxxlabel{C++17}\ \begin{lstlisting}
/* C++17 风格代码示例 */
#include <optional>
#include <variant>
#include <string_view>
#include <filesystem>
#include <vector>
#include <map>

struct Point { int x, y; };

std::optional<int> find_value(const std::map<std::string, int>& m,
                              std::string_view key) {
    if (auto it = m.find(std::string(key)); it != m.end()) {
        return it->second;            // 返回 optional
    }
    return std::nullopt;
}

template <typename T>
constexpr T square(T x) {
    if constexpr (std::is_integral_v<T>) {     // 编译期分支
        return x * x;
    } else {
        return x * x;
    }
}

int main() {
    std::map<std::string, int> ages{{"Alice", 30}, {"Bob", 25}};

    // 结构化绑定 + optional
    if (auto age = find_value(ages, "Alice")) {
        std::cout << "Alice is " << *age << "\n";
    }

    // 结构化绑定遍历 map
    for (const auto& [name, age] : ages) {
        std::cout << name << ": " << age << "\n";
    }

    // filesystem
    namespace fs = std::filesystem;
    for (const auto& entry : fs::directory_iterator(".")) {
        std::cout << entry.path() << "\n";
    }

    // variant
    std::variant<int, double, std::string> v = "hello";
    std::visit([](auto&& arg) { std::cout << arg << "\n"; }, v);

    return 0;
}
\end{lstlisting}

\section{C++20：另一次革命}

\yrlbl{2020} ISO/IEC 14882:2020，俗称 \textbf{C++20}。C++20 是 C++11 之后最大的一次更新，引入了四项"重型特性"，被合称为\textbf{C++20 四巨头}：

\begin{itemize}
  \item \textbf{Concepts}（概念）：\kw{template<typename T> requires Sortable<T>}，对模板参数的具名约束，替代 SFINAE 技巧。
  \item \textbf{Ranges}（范围）：\kw{views::filter}、\kw{views::transform}、\kw{views::take}，可组合的惰性视图，重塑 STL 算法。
  \item \textbf{Coroutines}（协程）：\kw{co\_await}、\kw{co\_yield}、\kw{co\_return}，语言层面的协程支持。
  \item \textbf{Modules}（模块）：\kw{import}、\kw{export}，替代头文件，编译模型根本变化。
\end{itemize}

C++20 的其他重要特性：

\begin{itemize}
  \item \textbf{Lambda 改进}：可捕获 \kw{*this}，支持模板参数。
  \item \textbf{\kw{consteval}}：强制编译期求值的函数（"立即函数"）。
  \item \textbf{\kw{constinit}}：强制静态变量编译期初始化。
  \item \textbf{三路比较 \kw{<=>}}：宇宙飞船运算符，一次性定义全部比较。
  \item \textbf{\kw{std::span}}：非拥有的连续序列视图。
  \item \textbf{\kw{std::format}}：类型安全的格式化（来自 Python/Rust 风格）。
  \item \textbf{\kw{std::jthread}}：自动 join 的线程。
  \item \textbf{\kw{std::stop\_token}}：协作式线程取消。
  \item \textbf{\kw{std::atomic} 增强}：\kw{atomic<shared\_ptr>}、\kw{atomic\_ref}、\kw{wait/notify}。
  \item \textbf{\kw{std::bit\_cast}}：类型双关（type punning）的标准方式。
  \item \textbf{\kw{<numbers>} 头文件}：\kw{std::numbers::pi} 等数学常量。
  \item \textbf{\kw{<source\_location>}}：替代 \kw{\_\_FILE\_\_}、\kw{\_\_LINE\_\_}。
  \item \textbf{\kw{<syncstream>}}：线程安全的输出流。
  \item \textbf{\kw{[[likely]]} / \kw{[[unlikely]]} 属性}：分支预测提示。
  \item \textbf{\kw{[[no\_unique\_address]]}}：空成员优化。
  \item \textbf{指定初始化器}：\kw{Point p\{.x = 1, .y = 2\};}（与 C99 同步）。
  \item \textbf{基于范围的 \kw{for} 改进}：\kw{for (init; auto\& x : container)}。
  \item \textbf{\kw{char8\_t} 类型}：UTF-8 字符的独立类型。
\end{itemize}

\cppxxlabel{C++20}\ \begin{lstlisting}
/* C++20 风格代码示例 */
#include <concepts>
#include <ranges>
#include <vector>
#include <format>
#include <coroutine>

// Concepts：约束模板参数
template <typename T>
concept Numeric = std::integral<T> || std::floating_point<T>;

template <Numeric T>
T add(T a, T b) { return a + b; }

// Ranges：可组合视图
void process(const std::vector<int>& data) {
    auto view = data
        | std::views::filter([](int x) { return x % 2 == 0; })
        | std::views::transform([](int x) { return x * x; })
        | std::views::take(5);

    for (int x : view) {
        std::print("{}\n", x);    // C++23 的 std::print
    }
}

// 三路比较
struct Version {
    int major, minor, patch;
    auto operator<=>(const Version&) const = default;
};

// Coroutines（生成器风格）
struct Generator {
    struct promise_type { /* ... */ };
    // 简化示意，完整实现需要 promise_type 与 iterator
};

Generator sequence() {
    for (int i = 0; i < 10; ++i) {
        co_yield i * i;
    }
}

int main() {
    auto result = add(1, 2);            // OK：int 满足 Numeric
    // auto bad = add(std::string("a"), std::string("b"));  // 编译错误

    Version v1{1, 2, 3}, v2{1, 3, 0};
    bool less = v1 < v2;                // 自动从 <=> 生成 < <= > >= == !=

    std::vector<int> data{1, 2, 3, 4, 5, 6, 7, 8, 9, 10};
    process(data);

    return 0;
}
\end{lstlisting}

\section{C++23：完善与同步}

\yrlbl{2023} ISO/IEC 14882:2023 发布，俗称 \textbf{C++23}。C++23 是 C++20 之后的完善版本，引入若干重要特性：

\begin{itemize}
  \item \textbf{\kw{std::print} 与 \kw{std::println}}：格式化输出（来自 \kw{std::format}）。
  \item \textbf{\kw{std::mdspan}}：多维数组视图（科学计算）。
  \item \textbf{\kw{std::expected<T, E>}}：错误处理的替代方案。
  \item \textbf{\kw{std::flat\_map}、\kw{std::flat\_set}}：基于排序数组的关联容器。
  \item \textbf{\kw{std::generator}}：协程生成器的标准实现。
  \item \textbf{\kw{std::stacktrace}}：栈追踪。
  \item \textbf{\kw{if consteval}}：编译期分支判断。
  \item \textbf{\kw{static operator()} 与 \kw{static operator[]}}：无捕获 lambda 优化。
  \item \textbf{\kw{deducing this}（P0847R7）}：成员函数的 \kw{this} 参数显式化，统一处理 const/非 const/左值/右值。
  \item \textbf{\kw{std::to\_underlying}}：枚举转底层类型。
  \item \textbf{多维下标 \kw{operator[]}}：\kw{a[i, j, k]}。
  \item \textbf{\kw{<stdbit.h>}}：位操作工具。
  \item \textbf{\kw{<latch>}、\kw{<barrier>}、\kw{<semaphore>}}：并发原语补全。
  \item \textbf{\kw{std::atomic<T>::wait} / \kw{notify}}：原子等待。
  \item \textbf{\kw{<span>} 增强}：\kw{std::subspan}。
  \item \textbf{移除}：\kw{std::experimental::} 命名空间下的部分组件。
\end{itemize}

\section{C++26：未来方向}

\yrlbl{2026} C++26 预计在 2026 年发布（截至本手册撰写时仍在制定中）。已经"基本确定"进入 C++26 的特性：

\begin{itemize}
  \item \textbf{Static Reflection}（静态反射，P2996）：编译期获取类型信息（成员名、类型、属性等）。这是 C++ 社区讨论了 15 年以上的特性，终于在 C++26 中通过。
  \item \textbf{Contracts}（契约，P2900）：函数前置/后置条件、断言。语法：\kw{void f(int x) [[pre: x > 0]] [[post r: r >= 0]];}。
  \item \textbf{Networking}（\texttt{<net>}、\texttt{<async>}）：异步 I/O、socket、执行器。基于 Boost.Asio 的设计。
  \item \textbf{\kw{std::execution}}：通用执行框架（"senders/receivers"）。
  \item \textbf{\kw{std::inplace\_vector}}：固定容量、堆栈分配的动态大小容器。
  \item \textbf{Pack indexing}：\kw{Args...[0]}，直接访问变参包中的特定参数。
  \item \textbf{\kw{constexpr} 异常}：编译期函数可以抛出与捕获异常。
  \item \textbf{\kw{<linear\_algebra>}}：基础线性代数库（提案中）。
  \item \textbf{Safety Profiles}（仍在讨论）：编译器辅助的安全检查（如边界、生命周期）。
\end{itemize}

C++26 的工作量与争议性都很大。Static Reflection 与 Contracts 各自经历了多版提案、数百页讨论、多次投票否决后重新设计。这是 ISO/IEC 治理流程的典型表现：\textbf{重大特性的通过往往需要 5-10 年的迭代}。

\begin{timelinebox}
\textbf{C++ 标准化时间线（1979-2026）}

\yrlbl{1979} Stroustrup 在 Bell Labs 启动"C with Classes"项目。

\yrlbl{1983} 语言定名为 C++，加入虚函数、运算符重载、引用。

\yrlbl{1985} AT\&T 发布 CFront 1.0，TC++PL 第一版出版。

\yrlbl{1989} 模板提案，C++ 2.0 加入异常、RTTI、多重继承、抽象类。

\yrlbl{1990} ARM（Annotated C++ Reference Manual）出版。

\yrlbl{1991} ANSI X3J16 委员会成立。

\yrlbl{1993} ISO/IEC WG21 委员会成立，与 X3J16 联合工作。

\yrlbl{1994} STL 被投票纳入标准库。

\yrlbl{1998} ISO/IEC 14882:1998（C++98）发布。

\yrlbl{2003} ISO/IEC 14882:2003（C++03，技术勘误）发布。

\yrlbl{2003-2011} "C++0x 休眠期"，委员会缓慢推进设计讨论。

\yrlbl{2011} ISO/IEC 14882:2011（C++11）发布——"一门新语言"。

\yrlbl{2014} ISO/IEC 14882:2014（C++14）发布。

\yrlbl{2017} ISO/IEC 14882:2017（C++17）发布。

\yrlbl{2020} ISO/IEC 14882:2020（C++20）发布——"C++20 是新的 C++11"。

\yrlbl{2023} ISO/IEC 14882:2023（C++23）发布。

\yrlbl{2026} C++26 预计发布（Reflection、Contracts、Networking）。
\end{timelinebox}

% ====================================================================
\part{ISO/IEC 治理结构}
% ====================================================================

\chapter{ISO/IEC SC22 与 WG14/WG21}

\section{ISO 与 IEC 简介}

\govlabel\ \textbf{ISO}（International Organization for Standardization，国际标准化组织）成立于 1947 年，总部设在日内瓦。其成员是\textbf{各国国家标准机构}（National Body，NB），如美国的 ANSI、英国的 BSI、德国的 DIN、法国的 AFNOR、日本的 JISC、中国的 SAC（国家标准化管理委员会）。ISO 不直接接受公司或个人会员——任何想参与 ISO 工作的实体都必须通过本国 NB。

\textbf{IEC}（International Electrotechnical Commission，国际电工委员会）成立于 1906 年，负责电气与电子领域的标准。信息技术（IT）领域同时涉及 ISO 与 IEC 的范畴，因此 1987 年成立了 \textbf{ISO/IEC JTC1}（Joint Technical Committee 1，联合技术委员会 1），统一管理 IT 标准。

\textbf{ISO/IEC JTC1} 下设多个 \textbf{SC}（Sub-Committee，分技术委员会），其中：

\begin{itemize}
  \item \textbf{SC22}：编程语言、环境及其软件接口（Programming languages, their environments and system software interfaces）。C、C++、Ada、Fortran、POSIX 等都归 SC22 管理。
  \item \textbf{SC32}：数据管理与交换（Data management and interchange）。SQL 归 SC32 管理。
  \item \textbf{SC29}：音频、图像、多媒体与超媒体信息编码。JPEG、MPEG 归 SC29 管理。
  \item \textbf{SC27}：信息安全、网络安全与隐私保护。
  \item \textbf{SC38}：云计算与分布式计算平台。
  \item \textbf{SC41}：物联网（IoT）与数字孪生。
  \item \textbf{SC42}：人工智能。
\end{itemize}

\textbf{SC22} 下设多个 \textbf{WG}（Working Group，工作组），按语言/标准划分：

\begin{itemize}
  \item \textbf{WG14}：C 语言（Programming language C），负责 ISO/IEC 9899。
  \item \textbf{WG21}：C++ 语言（Programming language C++），负责 ISO/IEC 14882。
  \item \textbf{WG5}：Fortran。
  \item \textbf{WG9}：Ada。
  \item \textbf{WG11}：POSIX（与 IEEE 1003 联合）。
  \item \textbf{WG15}：Ada 运行环境。
  \item \textbf{WG4}：COBOL（已休眠）。
  \item \textbf{WG20}：CJK 字符集国际化（与 Unicode 协调）。
\end{itemize}

C 与 C++ 标准由 WG14 与 WG21 分别负责，但两组同属 SC22，每年至少有一次联合会议（通常是夏季会议），讨论 C/C++ 之间的接口一致性问题。

\section{国家成员体（National Body，NB）的角色}

ISO/IEC 的所有正式参与都通过\textbf{国家成员体}进行。每个 ISO 成员国（约 170 个）有且仅有一个 NB，由该国国家标准机构担任。NB 在 ISO/IEC 中有两种身份：

\begin{itemize}
  \item \textbf{P 成员}（Participating member，参与成员）：有义务参与技术工作、对草案投票、可派出代表参加 WG 会议、可提交提案。WG14 与 WG21 的 P 成员通常包括：美国（ANSI/INCITS）、英国（BSI）、德国（DIN）、法国（AFNOR）、加拿大（SCC）、日本（JISC）、中国（SAC）、韩国（KATS）、瑞士（SNV）、荷兰（NEN）、北欧诸国等。
  \item \textbf{O 成员}（Observer member，观察员）：可获取标准草案、可观察会议进展，但\textbf{不投票}、不强制参与。
\end{itemize}

P 成员的投票是标准通过的关键。每份标准草案在 CD、DIS、FDIS 三个阶段都需要 P 成员投票，达到规定的多数才能进入下一阶段。

\section{NB 如何组织内部参与}

每个 NB 在自己国家内部组织"参与机制"。模式因国家而异：

\begin{itemize}
  \item \textbf{美国 ANSI/INCITS}：通过 INCITS（InterNational Committee for Information Technology Standards）下的技术委员会组织。WG14 与 WG21 的美国代表由 \textbf{INCITS PL22.11}（C）与 \textbf{PL22.16}（C++）委员会选出。任何美国组织（公司、大学、个人）都可以加入这些委员会，付费成为成员后获得投票权与提案权。
  \item \textbf{英国 BSI}：通过 IST/5（C）与 IST/5/-/1（C++）等小组组织。
  \item \textbf{德国 DIN}：通过 NIA-01-37 等小组组织。
  \item \textbf{日本 JISC}：通过 JIS X 3010（C）与 JIS X 3014（C++）的对应委员会组织。
  \item \textbf{中国 SAC}：通过 SAC/TC28（信息技术标准化技术委员会）下属的语言分组组织。
\end{itemize}

\begin{govbox}
\textbf{治理观察：国家代表与公司利益的张力}

ISO 的设计原则是"国家代表国家利益"，但现实中：

(1) WG21 美国代表团（PL22.16）的成员绝大多数来自美国公司——Microsoft、Google、Apple、Facebook、IBM、Intel、NVIDIA、Bloomberg 等。他们的投票立场往往反映雇主的技术偏好，而不是抽象的"美国国家利益"。

(2) 同一公司在多个 NB 中都有代表（Microsoft 在 ANSI、BSI、DIN 都有派员），可以协调投票立场。

(3) 个人也可以加入 ANSI/INCITS 委员会获得投票权——这是 ISO 体系相对开放的一面。WG21 中有相当比例的"独立专家"，他们以个人身份参与，不代表任何公司。

这种"国家代表 + 公司利益 + 独立专家"的混合模式，使得 ISO 治理结构既不完全去中心化（公司利益仍然存在），也不中心化（没有任何单一公司能控制投票）。
\end{govbox}

\section{WG14 与 WG21 的规模与组成}

\textbf{WG14}（C）相对较小：

\begin{itemize}
  \item P 成员约 15-20 个国家。
  \item 每次现场会议约 30-50 名代表。
  \item 主席（Convener）历史上长期由 David Keaton 担任，目前由 Jean Heyd Meneide 等共同主持。
  \item 工作节奏：每年 2-3 次会议，每次 1 周。
\end{itemize}

\textbf{WG21}（C++）规模较大：

\begin{itemize}
  \item P 成员约 20-25 个国家。
  \item 每次现场会议约 100-200 名代表。
  \item 主席（Convener）历史上由 Bill Gibbons、Tom Plum 等担任，目前是 Thomas Köppe（Google）。
  \item 工作节奏：每年 3 次会议（冬季 2 月、夏季 6-7 月、秋季 11 月），每次 1 周；外加每周的远程电话会议（telecons）。
  \item WG21 内部细分为多个子小组（CWG、LWG、EWG、LEWG、SG），各自有自己的邮件列表与会议时段。
\end{itemize}

\chapter{提案、投票与标准发布流程}

\section{提案的来源与提交}

任何 WG21 P 成员（包括公司代表、国家代表、独立专家）都可以提交\textbf{提案}（proposal）。提案以\textbf{编号文档}的形式发布在 WG21 官方网站 \seqsplit{open-std.org/jtc1/sc22/wg21/docs/papers/}：

\begin{itemize}
  \item \textbf{N 编号}：早期文档（N1000-N5000），如 \paper{N3225}（移动语义）、\paper{N2668}（lambda）。
  \item \textbf{P 编号}：2014 年起的新编号系统，\paper{P0732R2}（结构化 NTTP）、\paper{P1907R1}（浮点 NTTP）、\paper{P2996}（反射）。R 后数字是修订版本号（R0 初版，R1 第一修订...）。
  \item \textbf{WG21 内部文档}：CWG/LWG 的 issues list、reflectors 邮件等。
\end{itemize}

提案的标准结构：

\begin{enumerate}
  \item \textbf{动机}（Motivation）：为什么需要这个特性。
  \item \textbf{设计}（Proposal/Wording）：具体的语法与语义。
  \item \textbf{影响}（Impact）：对现有代码、编译器、标准库的影响。
  \item \textbf{实现}（Implementation Experience）：哪些编译器/库已经实现，性能与正确性数据。
  \item \textbf{替代方案}（Alternatives Considered）：讨论被否决的设计。
  \item \textbf{规范文本}（Proposed Wording）：直接插入标准草案的精确措辞。
\end{enumerate}

WG21 对"实现经验"的要求越来越高。一个未经任何编译器实现的提案，几乎不可能通过——委员会要求至少 2-3 个主流编译器（GCC、Clang、MSVC）有实验性实现，并能展示真实代码的运行效果。

\section{讨论与共识形成}

提案提交后经历以下讨论流程：

\begin{enumerate}
  \item \textbf{邮件列表（reflector）讨论}：提案作者发到 WG21 reflector（isocpp.org 的 std-proposals 邮件列表），社区成员评论、提问、批评。这一步可能持续数月。
  \item \textbf{EWG/LEWG 早期评审}：根据提案性质，由 \textbf{EWG}（Evolution Working Group，语言演化）或 \textbf{LEWG}（Library Evolution Working Group，库演化）评审。评审重点：方向是否正确、设计是否合理、是否值得继续投入。
  \item \textbf{CWG/LWG 措辞评审}：方向通过后，由 \textbf{CWG}（Core Working Group，语言核心）或 \textbf{LWG}（Library Working Group，库）做精确措辞评审。这一步极其细致，会逐字逐句推敲标准文本。
  \item \textbf{Straw Poll}（不记名民意调查）：在评审过程中，主席会发起 straw poll，让代表表达倾向。Straw poll 不具约束力，但能反映共识强度。常见的 straw poll 形式："支持进入下一阶段的人数？反对？弃权？"
  \item \textbf{Plenary 投票}：提案在 CWG/LWG 完成后，提交到 WG21 全体会议（plenary）正式投票。投票方式：举手或电子投票，P 成员代表按国家集体投票（每个国家一票，由代表团团长决定立场）。
\end{enumerate}

\begin{govbox}
\textbf{治理观察：共识 vs 多数决}

WG21 的工作文化强调"\textbf{共识决策}"（consensus），而不是简单多数决。理想流程是：

(1) 提案经过多轮讨论，作者根据反馈修订。

(2) Straw poll 显示明显支持（如 80\%+）时，提案进入下一阶段。

(3) 如果有强烈反对（哪怕只有少数人），提案会被退回继续修改，而不是强行投票。

(4) 最终在 plenary 投票时，几乎不会出现 51:49 这种胶着——因为前面已经形成了广泛共识。

这种"慢但稳"的决策模式有得有失：
\begin{itemize}
  \item \textbf{好处}：避免草率决策，保证标准的技术质量；让少数派意见被认真对待。
  \item \textbf{坏处}：进程极慢，一个有争议的特性可能讨论 5-10 年；某些"明显好"的提案因少数人反对而被拖延。
\end{itemize}
\end{govbox}

\section{NB 投票与标准阶段}

提案通过 plenary 投票后，进入\textbf{标准发布流程}。ISO 把标准从草案到正式发布分为若干阶段，每个阶段都有正式的 NB 投票：

\begin{enumerate}
  \item \textbf{NWIP}（New Work Item Proposal，新工作项目提案）：决定启动一项新标准。投票门槛：5 个 P 成员同意 + 多数 P 成员赞成。
  \item \textbf{WD}（Working Draft，工作草案）：委员会内部草稿，可能有多版。WD 阶段不进行 NB 投票，只在 WG 内部讨论。
  \item \textbf{CD}（Committee Draft，委员会草案）：第一次正式 NB 投票。所有 P 成员对 CD 内容投票（赞成/反对/弃权），同时提交 comments。投票期通常 2-3 个月。如果有重大反对意见，CD 会被修订并重新投票（CD2、CD3...）。
  \item \textbf{DIS}（Draft International Standard，国际标准草案）：第二次 NB 投票，向所有 ISO 成员（包括 O 成员）开放。投票期 5 个月。通过门槛：\textbf{2/3 以上 P 成员赞成} + \textbf{3/4 以上总票数赞成}（含 O 成员）。
  \item \textbf{FDIS}（Final Draft International Standard，国际标准最终草案）：最后一次 NB 投票。投票期 2 个月。通过门槛：\textbf{2/3 以上 P 成员赞成} + \textbf{3/4 以上总票数赞成}。FDIS 阶段不允许实质性修改，只能做编辑性修订。
  \item \textbf{IS}（International Standard，国际标准）：正式发布。ISO/IEC 联合签署，标准编号确定（如 ISO/IEC 14882:2020）。
\end{enumerate}

\begin{timelinebox}
\textbf{C++20 的发布流程（实例）}

\yrlbl{2018-03} Jacksonville 会议：C++20 进入 CD 阶段。

\yrlbl{2019-07} Cologne 会议：CD2 投票，进入 DIS 阶段。

\yrlbl{2019-07 至 2019-12} DIS 投票期，所有 ISO 成员可投票。

\yrlbl{2020-02} Prague 会议：DIS 投票通过，进入 FDIS 阶段。

\yrlbl{2020-09} FDIS 投票结束，赞成票 22:0（无反对）。

\yrlbl{2020-12-15} ISO/IEC 14882:2020 正式发布。

整个流程耗时约 30 个月，符合 ISO 标准发布的典型时间。
\end{timelinebox}

\section{缺陷报告（Defect Report, DR）}

标准发布后，使用者会发现规范中的歧义、矛盾、错误。这些问题通过\textbf{缺陷报告}（DR）流程处理：

\begin{enumerate}
  \item 任何人（不限于 P 成员）都可以向 CWG/LWG 提交 issue。
  \item 工作组评审 issue，分类为：NAD（Not A Defect，非缺陷）、CD3（延期到下版）、TC1（纳入下次技术勘误）、Defect（确认是缺陷，立即修正）。
  \item 确认的缺陷被纳入"标准修订"——通过 Technical Corrigendum（TC，技术勘误）发布。TC 不需要完整的 DIS/FDIS 流程，只需 WG 投票通过即可。
  \item 缺陷修正后的标准会标注修订版本，如 ISO/IEC 14882:2017/Cor 1:2018。
\end{enumerate}

DR 流程让标准能持续维护。C++ 标准自 C++98 以来已经处理了数千个 DR——这个数字反映了"用自然语言精确描述编程语言语义"的固有难度。

\chapter{WG21 的内部结构与文化}

\section{子工作组的划分}

WG21 内部按主题划分为多个子组，每个子组有自己的主席、邮件列表、会议时段：

\begin{itemize}
  \item \textbf{CWG}（Core Working Group，核心语言工作组）：处理语言核心问题——语法、语义、类型系统、模板、表达式、声明。主席历史上由 Richard Smith（Google，Clang 维护者）担任。CWG 维护"core issues list"，约 300+ 个开放 issue。
  \item \textbf{LWG}（Library Working Group，库工作组）：处理标准库问题——容器、算法、I/O、字符串、迭代器、分配器等。维护"library issues list"。
  \item \textbf{EWG}（Evolution Working Group，演化工作组）：处理新特性的设计方向——讨论提案是否值得纳入、设计是否合理。EWG 的产出是"方向决定"，不写标准文本。
  \item \textbf{LEWG}（Library Evolution Working Group，库演化工作组）：与 EWG 平行，但聚焦库的设计方向。
  \item \textbf{SG}（Study Group，研究组）：针对特定主题的研究组，每个 SG 有专门的领域：
    \begin{itemize}
      \item SG1：并发（Concurrency）。
      \item SG2：模块（Modules）。
      \item SG3：编译期计算（Compile-time / constexpr）。
      \item SG4：安全（Safety，包括边界、生命周期）。
      \item SG5：交易内存（Transactional Memory，已休眠）。
      \item SG6：数值计算（Numerics）。
      \item SG7：反射（Reflection）。
      \item SG8：网络（Networking）。
      \item SG9：范围（Ranges）。
      \item SG10：源代码能力探测（Source code capabilities，\kw{\_\_has\_feature} 等）。
      \item SG12：未定义行为（Undefined Behavior）。
      \item SG13：源代码安全（Source code security，已并入 SG4）。
      \item SG14：低延迟（Low Latency，游戏、HFT、嵌入式）。
      \item SG15：工具链（Tool Dependence，模块、构建系统）。
      \item SG16：Unicode。
      \item SG17：模式与方言（Patterns and dialects）。
      \item SG18：词法（Lexing）。
      \item SG19：机器学习（Machine Learning）。
      \item SG20：教育（Education）。
      \item SG21：契约（Contracts）。
      \item SG22：C/C++ 协调（C/C++ Liaison，与 WG14 联合）。
      \item SG23：安全规范（Safety specification）。
    \end{itemize}
\end{itemize}

\section{会议结构与节奏}

WG21 每年召开 3 次正式会议（春季 2 月、夏季 6-7 月、秋季 11 月），每次约 1 周，由不同国家轮办：

\begin{itemize}
  \item \textbf{周一到周五}：每天 9:00-17:00 工作。上午通常是各 SG/CWG/LWG/EWG/LEWG 的并行会议；下午有 plenary（全体会议）。
  \item \textbf{并行会议}：同一时段有多个小组开会，代表选择自己感兴趣的小组参加。比如 CWG 在 Room A 讨论 issue，EWG 在 Room B 评审提案，SG14 在 Room C 讨论低延迟特性。
  \item \textbf{Plenary 投票}：通常在周三与周五，所有代表集中投票当天评审过的提案。
  \item \textbf{晚餐与社交}：会议期间的非正式交流对共识形成非常重要——很多争议是在晚餐桌上解决的。
\end{itemize}

疫情之后，WG21 引入了\textbf{混合模式}（线下 + 线上）与\textbf{远程会议}（telecons）：

\begin{itemize}
  \item 每周 1-2 次 telecons，每次 2-3 小时，专门讨论某一类提案。
  \item 现场会议规模略减（部分代表远程参加）。
  \item Zoom/Slack 成为主要协作工具。
\end{itemize}

\section{委员会文化的特征}

WG21 的工作文化有几个鲜明特征：

\begin{itemize}
  \item \textbf{技术驱动}：讨论以技术论据为主，不太看"谁的级别高"。一位独立专家如果有充分的技术论据，可以与大公司代表平等讨论。
  \item \textbf{邮件优先}：所有重要讨论都发生在 reflector 上，现场会议主要是确认邮件中已经形成的共识。"没在邮件上讨论过的提案，会议上不会通过"是惯例。
  \item \textbf{措辞精确}：CWG/LWG 对标准文本的措辞要求极高——每个词、每个标点都被反复推敲。这是因为标准文本将来要作为编译器实现的依据，任何歧义都会引发争议。
  \item \textbf{向后兼容至上}：C++ 标准委员会有一条铁律——"不破坏现有代码"。任何会让现有合法 C++ 程序变非法或不行为变化的提案，几乎必然被否决。
  \item \textbf{实现验证}：纯理论提案不会通过，必须有编译器实现作为验证。
  \item \textbf{透明度}：所有提案文档、会议纪要、投票结果都公开发布在 open-std.org，任何人可查阅。
\end{itemize}

\begin{govbox}
\textbf{治理观察：委员会与编译器厂商的共生关系}

WG21 的成员大量来自编译器厂商：

\begin{itemize}
  \item \textbf{GCC}：Jason Merrill、Jonathan Wakely、Patrick Palka 等长期参与。
  \item \textbf{Clang/LLVM}：Richard Smith（CWG 主席）、Hubert Tong、Aaron Ballman 等。
  \item \textbf{MSVC}：Stephan T. Lavavej（STL 维护者）、Steve Downey 等。
  \item \textbf{Intel}：Anthony Williams（曾长期参与并发设计）。
  \item \textbf{IBM}、\textbf{Edison Design Group}（CFront 后继者）等。
\end{itemize}

这种"标准制定者 = 编译器实现者"的双重身份带来一个有趣的现象：委员会本身不生产编译器，但几乎所有标准实现都来自委员会成员所在的厂商。提案在通过前往往已经在 GCC/Clang/MSVC 中实现——这种"先实现再标准化"的模式让标准与现实保持一致。

但反面是：委员会没有自己控制的"参考实现"。如果某家公司退出（比如 Microsoft 关闭 MSVC），委员会无法强制其他厂商接手——这与 W3C、ECMA 等组织拥有 reference implementation 的模式不同。
\end{govbox}

% ====================================================================
\part{治理模式对比与客观陈述}
% ====================================================================

\chapter{三种治理模式：中心化、准中心化、去中心化}

\section{模式定义}

主流编程语言的标准/规范治理大致可分三类：

\begin{enumerate}
  \item \textbf{中心化治理}（Centralized governance）：语言规范与商标由\textbf{单一公司}（或单一个人）完全控制。公司单方面决定语言演化方向，规范文档可以不经任何外部投票发布。社区参与通常是咨询性的，没有否决权。

  \item \textbf{准中心化治理}（Quasi-centralized governance）：语言规范\textbf{名义上}归属一个开放组织（基金会、标准组织），但实际上由\textbf{单一公司主导}——公司控制了组织的多数投票权、商标、参考实现、生态基础设施。其他公司可参与讨论，但很难推翻主导公司的决定。

  \item \textbf{去中心化治理}（Decentralized governance）：语言规范由\textbf{多方独立组织共同参与}的标准机构（如 ISO/IEC、IEEE、Ecma）制定。决策需要多国/多公司共识，任何单一实体都不能控制结果。商标通常归属标准机构本身，而非任何公司。
\end{enumerate}

\section{代表性语言分类}

\begin{center}
\small
\begin{tabularx}{\textwidth}{@{}l l X@{}}
\toprule
\textbf{语言} & \textbf{治理模式} & \textbf{说明} \\
\midrule
Java & 中心化 & 商标与规范由 Oracle 控制，OpenJDK 是参考实现但 Oracle 决定发布。 \\
Kotlin & 中心化 & JetBrains 主导，Kotlin Foundation 名义上拥有商标但实际由 JetBrains 控制。 \\
Swift & 中心化（演化中） & Apple 主导；2018 开源后成立 Swift Core Team，但 Apple 仍占多数席位。 \\
Rust & 准中心化 & Rust Foundation 拥有商标，Mozilla 是创始成员；Rust Core Team 由原 Mozilla 团队成员主导。 \\
Go & 准中心化 & Go 由 Google 设计、开发、决策；Go Team 是 Google 内部组织，外部贡献者无投票权。 \\
C\# & 准中心化 & Microsoft 主导设计，ECMA-334 标准化，.NET Foundation 托管生态；Microsoft 仍是 .NET Foundation 最大成员。 \\
TypeScript & 准中心化 & Microsoft 完全控制，未通过任何标准组织。 \\
Python & 准中心化 & PSF 持有商标，2018 年前 Guido van Rossum 任 BDFL，之后改为 Steering Council（5 人，PSF 选举）。 \\
C & 去中心化 & ISO/IEC JTC1/SC22/WG14 制定，无单一公司控制。 \\
C++ & 去中心化 & ISO/IEC JTC1/SC22/WG21 制定，无单一公司控制。 \\
SQL & 去中心化 & ISO/IEC JTC1/SC32 制定。 \\
POSIX & 去中心化 & IEEE 1003 / ISO/IEC 9945 联合制定。 \\
ECMAScript & 去中心化 & Ecma International TC39 制定（Ecma 是欧洲标准组织，模式接近 ISO）。 \\
\bottomrule
\end{tabularx}
\end{center}

\section{关键判断维度}

判断治理模式可以从以下维度入手：

\begin{itemize}
  \item \textbf{商标持有者}：语言名称与 logo 归谁所有？（如"Java"商标归 Oracle，"C++"无任何公司拥有，"JavaScript"商标归 Oracle 但 ECMA 持有 ECMAScript 规范）
  \item \textbf{规范制定者}：谁有权发布新版本的规范？需要哪些步骤？
  \item \textbf{参考实现}：是否有官方实现？谁维护？源代码许可证是什么？
  \item \textbf{投票权分配}：谁有投票权？是按公司、按国家、按个人？是否需要多数同意？
  \item \textbf{外部参与门槛}：第三方（公司、个人）能否参与？需要什么资格？费用多少？
  \item \textbf{决策透明度}：会议、提案、投票记录是否公开？
  \item \textbf{退出机制}：如果主导者放弃项目，规范会怎样？
\end{itemize}

\begin{govbox}
\textbf{治理观察：模式不是非黑即白}

实际的治理模式往往是混合的：

\begin{itemize}
  \item \textbf{Java}：商标与规范决策中心化（Oracle），但 OpenJDK 是开源项目，社区贡献通过 JEP（JDK Enhancement Proposal）流程，部分外部参与者能影响设计。
  \item \textbf{C++}：规范决策去中心化（ISO/IEC），但实现严重依赖少数厂商（GCC、Clang、MSVC），如果这三家都拒绝实现某个特性，该特性等同于不存在。
  \item \textbf{Python}：商标归 PSF（基金会），技术上由 Steering Council 决策，但实际生态由少数大公司（Google、Meta、Microsoft）资助。
\end{itemize}

判断治理模式时应看\textbf{决策权}与\textbf{商标控制权}，而不是开源协议或参与形式。
\end{govbox}

\chapter{去中心化治理的优势：客观陈述}

\section{优势一：不受单一公司决策影响}

去中心化模式最根本的优势：\textbf{规范的演化方向不由任何一家公司单方面决定}。这与中心化模式形成对比。

\subsection{Java 商标战与 Oracle 的决策权}

\yrlbl{2010} Oracle 收购 Sun Microsystems，获得 Java 商标与 OpenJDK 控制权。收购后 Oracle 做出一系列单方面决定：

\begin{itemize}
  \item \yrlbl{2010-2017} Java 发布节奏由 Oracle 决定，长期保持"3-4 年一个大版本"的缓慢节奏，与社区期望不符。
  \item \yrlbl{2017} Oracle 改变 Java SE 的许可证，对企业生产使用收费。这一决定完全由 Oracle 单方面做出，没有任何外部投票。
  \item \yrlbl{2010-2016} Oracle 与 Google 就 Android 使用 Java API 进行长达 6 年的诉讼（最终上诉至美国最高法院）。Oracle 的诉讼策略反映其商标控制权，但这一法律纠纷对整个 Java 生态产生深远影响。
  \item \yrlbl{2014} Oracle 起诉 Google 案的一审判决期间，Java Community Process（JCP）的多个执行委员会成员公开表达对 Oracle 决策方式的不满，但 JCP 没有权力推翻 Oracle 的决定。
\end{itemize}

\yrlbl{2017} Oracle 宣布将 Java SE 转为 6 个月发布周期。这一转变受到社区欢迎，但\textbf{转变的决定权依然在 Oracle 手中}——社区无法主动推动这一变化，只能等待 Oracle 自己想通。

\subsection{C/C++ 没有"Oracle 时刻"}

对比之下，C/C++ 历史中没有出现任何一家公司能单方面改变标准走向的时刻。即使 Microsoft、Google、Apple 这些大公司在 WG21 中有大量代表，他们也必须：

\begin{itemize}
  \item 通过提案流程（提交 paper、邮件讨论、EWG/CWG 评审）。
  \item 在 straw poll 中获得多数支持。
  \item 在 NB 投票中获得 2/3 + 3/4 的多数。
\end{itemize}

如果某家公司试图推动一个对其有利但对生态有害的提案，其他 NB 与独立专家会反对，提案大概率被否决。这种"互相制衡"的结构是去中心化治理的核心特征。

\section{优势二：不受单一基金会政策影响}

准中心化模式下，语言名义上归基金会管理，但基金会政策可能被主导公司改变。

\subsection{.NET Foundation 与 C\#}

\yrlbl{2014} Microsoft 与外部合作伙伴共同创立 \textbf{.NET Foundation}，托管 .NET 生态（包括 C\#、Roslyn 编译器、ASP.NET、Entity Framework 等）。基金会董事会初期由 Microsoft 主导。

\yrlbl{2018} .NET Foundation 进行治理改革，引入"开放董事会选举"，让社区成员能投票选举部分董事。这一改革被宣传为"去 Microsoft 化"。

\yrlbl{2021} .NET Foundation 的董事会选举引发争议：当选董事中有几位是 Microsoft 员工或长期合作伙伴。社区批评选举过程被 Microsoft 影响力主导。最终部分董事辞职，基金会重新调整治理结构。

\yrlbl{2023-2024} .NET Foundation 经历了多轮管理层变动、预算缩减、项目退出。基金会的政策变化直接影响 C\# 生态（如某些项目的维护状态、社区资助计划）。但 C\# 语言本身的设计决策仍在 Microsoft 内部完成（C\# 设计团队是 Microsoft 部门），.NET Foundation 不直接参与语言设计。

\subsection{Rust Foundation 与 Mozilla}

\yrlbl{2020} Mozilla 大规模裁员，解散了 Rust 团队的大部分成员。Rust 项目的未来一度受到质疑。

\yrlbl{2021} Rust Foundation 成立，由 Mozilla、AWS、Google、Huawei、Microsoft 五家创始白金会员资助。Rust 商标与基础设施移交基金会。

\yrlbl{2021-至今} Rust Foundation 的董事会包含白金会员代表与项目代表（Rust Core Team 派出的代表）。但白金会员（5 家大公司）在董事会中占多数，理论上能影响基金会决策。Rust 语言设计仍由 Rust Core Team 决定（不受基金会控制），但\textbf{项目基础设施、商标、预算}由基金会管理。

\subsection{C/C++ 没有"基金会风险"}

C/C++ 标准由 ISO/IEC 直接管理。ISO/IEC 不是基金会，而是一个由 170 个国家成员体组成的国际条约组织。任何一家公司、一个国家都无法单方面退出 ISO 并影响 C/C++ 标准——即使所有大公司集体退出 WG21，标准本身依然存在（ISO/IEC 14882 是已发布的国际标准）。

\section{优势三：历史延续性与跨国法律地位}

ISO/IEC 标准具有以下特性：

\begin{itemize}
  \item \textbf{已发布的标准永久有效}。ISO/IEC 9899:1990（C90）虽然已被 C99/C11/C17/C23 取代，但作为历史标准依然可获取、可引用。
  \item \textbf{标准本身是国际法律实体}。在政府采购、国际合同、行业标准引用中，"符合 ISO/IEC 14882"具有法律效力，远超过"符合 Microsoft C\# 语言规范"。
  \item \textbf{标准文档可被任何国家采纳为本国标准}。中国的 GB/T 15273（C）、GB/T 18092（C++）就是直接采纳 ISO/IEC 9899/14882。
  \item \textbf{标准修订不需要重启组织}。即使 ISO 解散（不太可能），已发布的标准文本仍是历史记录。
\end{itemize}

\section{优势四：技术中立性}

去中心化治理要求决策必须考虑多方利益，这天然带来\textbf{技术中立性}：

\begin{itemize}
  \item 标准不会偏向任何特定平台（Windows、Linux、macOS）、特定硬件（x86、ARM）、特定厂商（Microsoft、Google、Apple）。
  \item 标准不会要求特定专利（ISO 有专利政策，要求专利持有人提供 RAND 许可，否则提案被否决）。
  \item 标准不会强制使用特定第三方库或服务。
\end{itemize}

例如 C++ 标准库的 \kw{std::filesystem} 必须在 Windows、Linux、macOS 上行为一致；C 标准的 \kw{<threads.h>} 必须能与 pthreads、Win32 threads、C11 原生线程都兼容实现。这种"跨平台中立"是任何一家公司主导的标准很难做到的（Microsoft 的 C\# 规范早期就有不少 Windows 特定假设）。

\section{优势五：退出风险低}

如果某家公司退出 ISO/IEC 工作（不再派代表、不再付费），影响有限：

\begin{itemize}
  \item 标准本身继续存在与演化。
  \item 其他公司与国家代表填补空缺。
  \item 该公司原有的提案、邮件讨论、issue 报告依然是公开记录。
\end{itemize}

历史上 AT\&T 退出 Bell Labs 的 C/C++ 工作、Sun Microsystems 被收购后的 Java 团队变动，都没有让标准本身受到致命打击。这种"反脆弱性"是去中心化治理的核心优势。

\begin{govbox}
\textbf{治理观察：去中心化优势的边界}

上述优势的\textbf{前提}是"标准已被广泛采纳"。如果一个 ISO 标准没有实现、没有用户，那么"不受单一公司控制"的优势就毫无意义——因为没有公司在乎这个标准。

C/C++ 之所以能享受去中心化治理的优势，是因为它们已经有 35+ 年的实现历史、巨大的用户基数、深厚的生态系统。一个新启动的语言选择 ISO 治理，可能反而会因为流程缓慢、缺乏厂商支持而失败。
\end{govbox}

\chapter{去中心化治理的代价：客观陈述}

\section{代价一：委员会不生产编译器}

这是去中心化治理最显著的代价。WG14/WG21 自身不维护任何编译器——\textbf{所有 C/C++ 编译器都由委员会成员所在的公司或志愿者实现}：

\begin{itemize}
  \item \textbf{GCC}：由 Free Software Foundation（FSF）维护，主要贡献者来自 Red Hat、Apple、SUSE、ARM 等。
  \item \textbf{Clang/LLVM}：由 LLVM Foundation 维护，主要贡献者来自 Apple、Google、Microsoft、Intel、AMD、ARM 等。
  \item \textbf{MSVC}：由 Microsoft 维护，闭源（部分组件开源），仅支持 Windows 平台。
  \item \textbf{ICC/ICX}（Intel）：基于 Clang 的 Intel 编译器。
  \item \textbf{其他}：Oracle Solaris Studio、IBM XL、Cray、PGI/NVIDIA HPC 等专有编译器。
\end{itemize}

\subsection{委员会与编译器厂商的协作模式}

WG21 不维护任何"参考实现"。这与 W3C 维护参考浏览器、Ecma 维护参考 JavaScript 引擎的模式不同。WG21 依赖\textbf{编译器厂商的实验性实现}来验证提案：

\begin{enumerate}
  \item 提案作者（通常也是某编译器厂商的工程师）实现一个原型。
  \item 原型被合并到主编译器的实验分支（如 Clang 的 \kw{-fexperimental-library}、GCC 的 \kw{-std=c++2b}）。
  \item 用户在真实代码中测试，反馈问题。
  \item 提案根据反馈修订，再次评审。
  \item 标准通过后，编译器实现转为正式支持。
\end{enumerate}

这种"标准依赖实现验证"的模式有得有失：

\begin{itemize}
  \item \textbf{好处}：标准不会脱离实际，每个特性都经过编译器与用户的检验。
  \item \textbf{坏处}：如果某家公司不愿投入资源实现，标准就会推迟。例如 \seqkw{std::thread} 在 MSVC 上迟到了好几年；C++17 的 \seqkw{std::filesystem} 在 GCC 与 Clang 上各自有不同的实验性命名空间；C++20 的 Modules 至今在三大编译器上的实现仍不完整。
\end{itemize}

\subsection{没有一致性测试套件}

ISO/IEC 不维护任何"一致性测试套件"（conformance test suite）。判断一个 C/C++ 编译器是否"符合标准"完全依赖：

\begin{itemize}
  \item 编译器厂商的自我声明（"我们支持 C++17"）。
  \item 第三方测试套件（如 Perennial 的 C 一致性测试、Plum Hall Validation Suite，这些都是商业产品）。
  \item 开源测试项目（如 csmith 用于发现 C 编译器 bug、cxxopts 用于测试 C++ 标准库一致性）。
  \item cppreference 与 godbolt 等社区资源。
\end{itemize}

对比之下，W3C 为 HTML/CSS 维护官方一致性测试套件（web-platform-tests），Ecma 为 ECMAScript 维护 test262，Java 有 TCK（Technology Compatibility Kit，由 Oracle 控制）。ISO/IEC 没有类似的官方机制——这反映了 ISO 的角色定位是"标准制定者"而非"实现监督者"。

\section{代价二：进程缓慢}

去中心化决策流程的代价是\textbf{极慢}。

\subsection{标准周期}

\begin{center}
\small
\begin{tabularx}{\textwidth}{@{}l l l X@{}}
\toprule
\textbf{标准} & \textbf{上一版} & \textbf{本版} & \textbf{间隔} \\
\midrule
C90 → C99 & 1990 & 1999 & 9 年 \\
C99 → C11 & 1999 & 2011 & 12 年 \\
C11 → C17 & 2011 & 2018 & 7 年（仅修订）\\
C17 → C23 & 2018 & 2024 & 6 年 \\
C++98 → C++03 & 1998 & 2003 & 5 年（仅修订）\\
C++03 → C++11 & 2003 & 2011 & 8 年 \\
C++11 → C++14 & 2011 & 2014 & 3 年 \\
C++14 → C++17 & 2014 & 2017 & 3 年 \\
C++17 → C++20 & 2017 & 2020 & 3 年 \\
C++20 → C++23 & 2020 & 2023 & 3 年 \\
C++23 → C++26 & 2023 & 2026（预计） & 3 年 \\
\bottomrule
\end{tabularx}
\end{center}

C++ 自 C++11 之后稳定在"3 年一版"的节奏，这已经是 ISO 体系下相对快的速度。对比 SQL 标准（约 5-8 年一版）、Ada（约 10 年一版）、POSIX（约 15 年一版），C++ 是 SC22 中迭代最快的标准之一。

\subsection{单个提案的周期}

单个特性的提案周期通常更长：

\begin{itemize}
  \item \textbf{Concepts}：首次提案 2003 年（C++0x 项目），C++20 才通过，历时 \textbf{17 年}。期间被否决并重新设计两次。
  \item \textbf{Modules}：首次提案 2005 年，C++20 通过，历时 \textbf{15 年}。
  \item \textbf{Coroutines}：首次提案 2004 年，C++20 通过，历时 \textbf{16 年}。
  \item \textbf{Static Reflection}：首次提案 2008 年，C++26 预计通过，历时 \textbf{18 年}。
  \item \textbf{Contracts}：首次提案 2003 年，C++26 预计通过，历时 \textbf{23 年}。
  \item \textbf{Ranges}：首次提案 2012 年，C++20 通过，历时 \textbf{8 年}。
  \item \textbf{\kw{std::format}}：首次提案 2017 年，C++20 通过，历时 \textbf{3 年}（相对快）。
\end{itemize}

这种"5-15 年的提案周期"对快速演化的领域（如 ML、Web）是难以接受的。Rust、Swift、Kotlin 等语言能以"年度发布"的节奏演化，正是因为它们不需要 ISO 流程。

\subsection{进程缓慢的具体表现}

提案被拖延的常见原因：

\begin{itemize}
  \item 设计争议：不同公司/专家有根本性分歧（如 Concepts 的"约束语法"经过 5+ 版重新设计）。
  \item 实现验证不足：没有足够多的编译器原型，无法证明特性可行。
  \item 措辞复杂性：CWG/LWG 对标准文本逐字推敲，一个 issue 可能讨论 1-2 年。
  \item 兼容性问题：发现提案会破坏现有代码，必须重新设计。
  \item 优先级冲突：委员会精力有限，提案被排队等待。
\end{itemize}

\section{代价三：投票政治化}

虽然 ISO/IEC 设计上是"国家代表国家利益"，实际运作中投票往往反映\textbf{公司与组织利益}：

\begin{itemize}
  \item 美国代表团（PL22.16）的成员主要是 Microsoft、Google、Apple、Meta、NVIDIA、Bloomberg、IBM 的员工，他们的投票立场往往与雇主的技术战略一致。
  \item 同一公司在多个 NB 中派员（如 Microsoft 在 ANSI、BSI、DIN 都有代表），可在不同投票中协调立场。
  \item 中小公司与独立专家的参与门槛较高：参加 ANSI 委员会需要年费（数千美元），现场会议差旅费高（每年 3 次国际旅行）。
\end{itemize}

这种"政治化"在重大设计决策上表现明显：

\begin{itemize}
  \item \yrlbl{2012} Concepts 第一版（"C++0x Concepts"）被否决——Microsoft 与 IBM 反对，Google 与 Facebook 支持。否决后 Concepts 重新设计为更简单的"C++20 Concepts"。
  \item \yrlbl{2018} Modules 设计争议：MSVC 团队、Apple Clang 团队、GCC 团队对 BMI（Binary Module Interface）格式有根本性分歧，最终标准只规定接口语义、不规定文件格式。
  \item \yrlbl{2020} Reflection 提案 P2996 在 EWG 评审中经历多次修订，作者团队（Bloomberg、Google、JetBrains）与反对者（部分 Clang 维护者）的辩论持续 3 年。
\end{itemize}

\section{代价四：标准文档收费}

ISO/IEC 标准的\textbf{正式文档是收费出版物}：

\begin{itemize}
  \item ISO/IEC 9899:2018（C17）：约 \textbf{218 瑞士法郎}（约 240 美元）一份 PDF。
  \item ISO/IEC 14882:2020（C++20）：约 \textbf{238 瑞士法郎}（约 260 美元）一份 PDF，1756 页。
  \item ISO/IEC 14882:2024（C++23）：约 \textbf{238 瑞士法郎}，1700+ 页。
  \item 每个版本单独销售，旧版本不免费提供。
\end{itemize}

ISO 的商业模式依赖标准销售收入。这与"开源精神"形成鲜明对比——Rust 规范、ECMAScript 规范、HTML 规范都免费提供。

\subsection{替代获取途径}

社区找到了几种免费获取 C++ 标准内容的方式：

\begin{itemize}
  \item \textbf{WG21 公开的工作草案}：open-std.org 上发布的 N4860（C++20 FDIS）、N4950（C++23 FDIS）等草案与正式标准在技术内容上完全相同，仅排版有差异。免费下载。
  \item \textbf{cppreference.com}：社区维护的 C++ 参考站点，按标准条目组织内容，免费且持续更新。
  \item \textbf{C++ Core Guidelines}：Stroustrup 与 Sutter 维护的"非官方但权威"的指南，免费。
\end{itemize}

这种"标准收费但实际免费可获取"的状况让 ISO 的商业模式略显尴尬——大多数实际用户不会购买官方文档，而是依赖工作草案与社区资源。

\section{代价五：国际组织的固有缺点}

ISO/IEC 作为国际条约组织，有一些固有缺点：

\subsection{官僚与流程冗长}

\begin{itemize}
  \item 标准发布流程必须经过 NWIP → WD → CD → DIS → FDIS → IS 六个阶段，每阶段都有固定时长（DIS 5 个月、FDIS 2 个月），无法压缩。
  \item 即使所有 P 成员都同意，也不能跳过任何阶段。
  \item ISO 秘书处的行政流程（编号、出版、翻译）需要额外时间。
\end{itemize}

\subsection{地域不平衡}

\begin{itemize}
  \item WG21 的活跃代表绝大多数来自北美与欧洲。亚洲（除日本外）、非洲、南美的参与度较低。
  \item 现场会议主要在欧美举办，亚洲代表参会成本高。
  \item 邮件讨论以英语为主，非英语母语者参与门槛高。
\end{itemize}

\subsection{主权代表 vs 技术专家的张力}

\begin{itemize}
  \item ISO 投票以"国家"为单位——一个国家一票。但每个国家代表团内部立场可能分歧（如美国代表团里 Microsoft 代表与 Google 代表立场不同），需要内部协调后再投出"美国一票"。
  \item 这种协调过程让技术决策受国家政治影响。例如某些标准在 DIS 投票时被某个国家整体反对（不是技术原因，而是政策原因）。
\end{itemize}

\subsection{无法强制执行}

\begin{itemize}
  \item ISO 不能要求任何公司实现标准。如果所有编译器厂商都不实现某个特性，该特性等同于不存在。
  \item ISO 不能审查编译器是否符合标准。"符合 ISO/IEC 14882"完全是厂商自我声明。
  \item ISO 不能阻止"标准分裂"——如果某家公司发布一个不兼容的扩展，ISO 无法干预。
\end{itemize}

\section{代价六："设计 by 委员会"风险}

委员会决策的固有缺点：\textbf{妥协产物可能不如单一设计师的设计清晰}。

\subsection{C++ 复杂度的争议}

C++ 标准文档已达 \textbf{1700+ 页}，远超 C（约 500 页）。这反映了 C++ 的特性数量与设计复杂度。常见批评：

\begin{itemize}
  \item 同一问题有多种解决方式（如智能指针有 \kw{unique\_ptr}、\kw{shared\_ptr}、\kw{weak\_ptr}，再加上传统裸指针）。
  \item 历史包袱重（如 \kw{std::auto\_ptr} 在 C++11 被废弃但保留到 C++17）。
  \item 概念之间的交互复杂（如模板、Concepts、SFINAE 之间的关系）。
\end{itemize}

部分批评者认为这是"设计 by 委员会"的必然结果——每个新特性都需要找到与现有特性的兼容方式，导致设计妥协。但支持者认为这种"渐进演化"恰恰保证了向后兼容与生态稳定。

\subsection{对比：单一设计师的语言}

\begin{itemize}
  \item \textbf{Go}：由 Robert Griesemer、Rob Pike、Ken Thompson 三人设计，Go Team（Google 内部）决定演化方向。Go 1.0 发布以来核心语法几乎没变（直到 Go 1.18 加入 Generics）。设计极其简洁。
  \item \textbf{Rust}：由 Graydon Hoare 个人发起，早期设计由 Graydon 与 Mozilla 团队完成，后续由 Rust Core Team 决定方向。Rust 的设计在"安全性"与"零开销抽象"之间保持高度一致。
  \item \textbf{Swift}：由 Chris Lattner 主导设计，Swift Core Team 决定演化。Swift 的设计在"现代性"与"易学性"之间相对一致。
  \item \textbf{Kotlin}：由 JetBrains 设计，JetBrains 完全控制演化方向。Kotlin 的设计简洁、聚焦实用性。
\end{itemize}

这些语言的演化速度都比 C++ 快得多，设计也更有"统一感"。但它们的代价是：演化方向高度依赖单一组织（Google、Mozilla/Rust Foundation、Apple、JetBrains）的战略决策。

\chapter{横向对比与客观陈述}

\section{多维度对比表}

\begin{center}
\scriptsize
\begin{tabularx}{\textwidth}{@{}p{2.2cm} p{2.6cm} p{2.6cm} p{2.6cm} p{2.6cm}@{}}
\toprule
\textbf{维度} & \textbf{中心化（Java/Oracle）} & \textbf{准中心化（C\#/.NET Foundation）} & \textbf{去中心化（C++/ISO/IEC）} & \textbf{单一个人主导} \\
\midrule
决策速度 & 快（公司单方面决定） & 中（公司主导，需协调生态） & 慢（多方共识，3 年周期） & 极快（个人决定） \\
\addlinespace
技术中立性 & 低（受公司战略影响） & 中（公司主导但有外部参与） & 高（多方互相制衡） & 低（受个人偏好影响） \\
\addlinespace
历史延续性 & 中（依赖公司存活） & 中（依赖基金会与公司） & 高（标准永久存在） & 低（个人退出即结束） \\
\addlinespace
实现保障 & 高（公司提供参考实现） & 高（公司主导实现） & 中（依赖厂商投入） & 低（无组织保障） \\
\addlinespace
用户控制权 & 低（用户无法投票） & 中（社区有部分参与） & 高（通过 NB 投票） & 极低（用户无影响） \\
\addlinespace
退出风险 & 高（公司可单方面放弃） & 中（基金会与公司双重依赖） & 低（标准独立存在） & 极高（个人即一切） \\
\addlinespace
商标控制 & 公司持有 & 基金会持有但公司主导 & ISO/IEC 持有 & 个人持有 \\
\addlinespace
文档获取 & 公司免费提供 & 公司免费提供 & ISO 收费但工作草案免费 & 个人决定 \\
\addlinespace
一致性测试 & 公司提供 TCK & 公司提供测试 & 无官方套件 & 无 \\
\addlinespace
设计一致性 & 高（公司战略一致） & 中（公司主导） & 中（委员会妥协） & 高（个人风格统一） \\
\bottomrule
\end{tabularx}
\end{center}

\section{各模式的代表性争议}

\subsection{Java：Oracle 时代的争议}

\begin{itemize}
  \item \yrlbl{2010-2016} Oracle vs Google 关于 Android 使用 Java API 的诉讼，最终上诉至美国最高法院（\yrlbl{2021} 判决 Google 胜诉，认定 API 使用属于 fair use）。
  \item \yrlbl{2017} Oracle 改变 Java SE 商业许可，对企业生产使用收费。社区分裂：部分用户转向 OpenJDK 发行版（AdoptOpenJDK、Corretto、Zulu）。
  \item \yrlbl{2019} Oracle 宣布 Java SE 8 的商业支持到期，企业必须升级到 Java SE 11+ 或迁移到 OpenJDK。这一决定推动大量企业迁移。
  \item \yrlbl{2020-至今} Java 发布节奏改为 6 个月一版（Oracle 单方面决定），社区对新版本质量与稳定性的批评增加。
\end{itemize}

\subsection{C\#/.NET Foundation：治理改革的反复}

\begin{itemize}
  \item \yrlbl{2014} Microsoft 创立 .NET Foundation，宣传"开放治理"。
  \item \yrlbl{2018} .NET Foundation 引入"开放董事会选举"，但实际选举结果仍被 Microsoft 影响力主导。
  \item \yrlbl{2021} 董事会选举争议，部分当选董事辞职，基金会重新调整治理。
  \item \yrlbl{2023-2024} 基金会预算缩减，多个项目退出，社区对基金会作用产生质疑。
  \item 整个过程中，C\# 语言设计仍由 Microsoft 内部 C\# 设计团队决定，.NET Foundation 不直接参与语言设计。
\end{itemize}

\subsection{Python：从 BDFL 到 Steering Council}

\begin{itemize}
  \item \yrlbl{1991-2018} Guido van Rossum 任 BDFL（Benevolent Dictator For Life，仁慈的终身独裁者），拥有最终决策权。这种"单一个人主导"模式让 Python 演化保持设计一致性，但也存在单点失效风险。
  \item \yrlbl{2018} Guido 因 PEP 572（赋值表达式 \kw{:=}）争议辞去 BDFL 职务。
  \item \yrlbl{2018-2019} Python 社区经过 PEP 8001-8016 系列提案讨论，最终选择"Steering Council"模式（5 人委员会，PSF 选举，2 年任期）。
  \item \yrlbl{2019-至今} Steering Council 模式下，Python 演化速度略有放缓，但避免了 BDFL 模式的单点风险。
\end{itemize}

\subsection{Rust：Mozilla 解散团队后的转折}

\begin{itemize}
  \item \yrlbl{2010-2020} Rust 由 Mozilla 资助开发，Rust Core Team 多数成员是 Mozilla 员工。
  \item \yrlbl{2020-08} Mozilla 大规模裁员，解散 Rust 团队大部分成员。Rust 项目未来受到质疑。
  \item \yrlbl{2021-02} Rust Foundation 成立，5 家白金会员（Mozilla、AWS、Google、Huawei、Microsoft）资助。商标与基础设施移交基金会。
  \item \yrlbl{2021-至今} Rust 演化未受 Mozilla 裁员影响，但治理结构从"准中心化（Mozilla 主导）"转为"准中心化（Rust Foundation 多方平衡）"。
\end{itemize}

\subsection{Swift：从 Apple 独家到核心团队的演化}

\begin{itemize}
  \item \yrlbl{2010-2014} Swift 由 Chris Lattner 主导在 Apple 内部秘密开发。
  \item \yrlbl{2014-06} WWDC 公开 Swift。
  \item \yrlbl{2015-12} Swift 开源，许可证 Apache 2.0。
  \item \yrlbl{2017} Swift.org 引入"Core Team"与"Evolution Process"，外部贡献者可通过 SE（Swift Evolution）提案参与设计。但 Core Team 多数成员仍是 Apple 员工。
  \item \yrlbl{2018} Chris Lattner 离开 Apple 加入 SiFive，后又加入 Google、Tesla、SiFive。Swift 的设计方向不再由单一设计师主导。
  \item \yrlbl{2020-至今} Swift 演化由 Swift Core Team 与 Evolution Workgroup 决定，Apple 仍是主导力量但不再是唯一决策者。
\end{itemize}

\section{客观陈述：哪种模式"更优"？}

本手册的立场是：\textbf{不评判哪种治理模式"更优"}——除非是\textbf{纯单个个人主导}的模式（如 BDFL）。

\subsection{为什么不下结论}

不同治理模式适合不同场景，没有绝对的优劣：

\begin{itemize}
  \item \textbf{快速演化的领域}（如 Web 前端、机器学习）适合中心化或准中心化模式。React、TensorFlow、PyTorch 的快速迭代依赖单一组织主导。
  \item \textbf{长期稳定的基础设施}（如操作系统接口、嵌入式语言）适合去中心化模式。POSIX、C、C++ 数十年不变的特性是去中心化治理的成果。
  \item \textbf{商业战略产品}（如 Java、C\#、Swift）适合中心化模式，公司能根据市场需求快速调整。
  \item \textbf{多方利益相关}（如 SQL、HTTP、ECMAScript）适合去中心化模式，避免任何一家独大。
\end{itemize}

\subsection{唯一的例外：纯单个个人主导}

本手册唯一明确的判断是：\textbf{纯单个个人主导}的治理模式具有公认的\textbf{单点失效风险}：

\begin{itemize}
  \item \textbf{Python BDFL 时代}：Guido 一旦健康出问题、辞职、去世，整个语言失去最终决策者。这种风险在 \yrlbl{2018} Guido 辞职时成为现实，迫使社区紧急设计新的治理结构。
  \item \textbf{Linux}：Linus Torvalds 个人拥有最终决策权，git 仓库的 maintainer 链条最终汇聚到他一人。Linus 的健康、性格、决策风格直接影响整个项目。社区曾多次讨论"Linus 之后"的问题，但没有正式答案。
  \item \textbf{Perl}：Larry Wall 个人主导设计，Perl 6（后改名 Raku）的长期开发延期与 Larry 的健康状况直接相关。
  \item \textbf{Ruby}：Matz（Yukihiro Matsumoto）拥有"最终决定权"，但实际通过 Ruby Association 与核心开发者团队运作。
\end{itemize}

这些案例的共同点：\textbf{一旦个人无法继续，整个语言的演化方向就失去锚点}。社区要么紧急重组治理（Python），要么陷入长期停滞（Perl 6），要么依赖个人长期投入（Linux）。

相比之下，组织主导的治理（无论中心化、准中心化、去中心化）都有"组织延续性"——即使个别领导者更替，组织本身仍能运作。这是为什么本手册认为"纯单个个人主导"是公认劣势的模式。

\section{C/C++ 选择的合理性}

回到 C/C++ 的具体情境：

\begin{itemize}
  \item \textbf{C 与 C++ 是基础设施级语言}。它们运行在操作系统、嵌入式设备、数据库、浏览器、游戏引擎、金融系统等关键领域。这些场景要求\textbf{稳定性}远超过\textbf{演化速度}。
  \item \textbf{多家公司同时投入实现}。GCC（FSF/Red Hat）、Clang（Apple/Google/Microsoft）、MSVC（Microsoft）、ICC（Intel）都在持续投入。没有任何一家公司能"接管"标准——这天然适合去中心化治理。
  \item \textbf{用户基数巨大且分散}。从嵌入式开发者到游戏工程师到金融量化分析师，C/C++ 用户横跨无数领域。任何单一公司都无法代表所有用户利益。
  \item \textbf{历史包袱深}。35+ 年的代码必须继续工作，向后兼容是铁律。这种约束让"快速演化"成为不可能，去中心化的"慢"反而成为优势。
\end{itemize}

对 C/C++ 来说，ISO/IEC 治理模式不是"最理想的选择"，而是"与生态现状最匹配的选择"。如果 C++ 在 1985 年选择由 AT\&T 独家控制，今天可能像 Java 一样面临 Oracle 收购后的争议；如果 C++ 选择由 Stroustrup 个人主导（像 BDFL），那么 2000 年代 Stroustrup 退休时可能像 Python 一样面临治理重组危机。

\section{未来展望}

C/C++ 标准治理的未来趋势：

\begin{itemize}
  \item \textbf{加速演化}：自 C++11 之后，C++ 稳定在 3 年一版的节奏，是 ISO 体系下相对快的速度。WG21 在保持流程的前提下努力提高效率（如使用 GitHub、Slack、Zoom 替代部分邮件）。
  \item \textbf{Safety 议题}：美国政府（CISA、NSA、白宫）近年多次发布"内存安全语言"倡议，将 C/C++ 列为"不安全"语言。这给 WG21 带来压力，需要在标准层面增加安全特性（如 Bounds Safety、Lifetime Safety）。SG4 与 SG23 正在讨论这些方向。
  \item \textbf{Reflection 与 Contracts}：C++26 的两项重磅特性将改变 C++ 编程范式。这些特性的通过过程体现了去中心化治理的"长期投入、最终共识"模式。
  \item \textbf{C 与 C++ 的协调}：WG14 与 WG21 通过 SG22 加强协调。C23 引入的 \kw{bool}、\kw{nullptr}、\kw{constexpr}、\kw{[[...]]} 等都与 C++ 同步。
  \item \textbf{与其他标准组织的合作}：与 IEEE（POSIX）、Ecma（ECMAScript）、W3C（Web）等组织在共享议题上协作（如 Unicode、网络、并发）。
\end{itemize}

\begin{govbox}
\textbf{治理观察：去中心化模式的边界条件}

本手册最后强调：去中心化治理的优势\textbf{只在特定条件下成立}：

(1) 标准必须已被广泛采纳，有真实的用户与实现。

(2) 必须有足够多的独立厂商参与（避免被某一两家垄断话语权）。

(3) 演化速度必须与场景匹配（基础设施语言可以慢，应用语言不行）。

(4) 必须有多方利益相关方互相制衡。

如果一个新语言选择 ISO/IEC 治理但缺乏这些条件，反而会因为流程缓慢、缺乏厂商支持而失败。Ada 的部分历史教训即在于此——Ada 标准化时美国政府是主要推动者，但商业化不足，最终市场份额被 C/C++ 取代。

C/C++ 的去中心化治理不是"放之四海皆准的最佳实践"，而是"与其生态现状匹配的一种选择"。
\end{govbox}

\end{document}

